%% 
%% Copyright 2019-2024 Elsevier Ltd
%% 
%% Version 2.4
%% 
%% This file is part of the 'CAS Bundle'.
%% --------------------------------------
%% 
%% It may be distributed under the conditions of the LaTeX Project Public
%% License, either version 1.2 of this license or (at your option) any
%% later version.  The latest version of this license is in
%%    http://www.latex-project.org/lppl.txt
%% and version 1.2 or later is part of all distributions of LaTeX
%% version 1999/12/01 or later.
%% 
%% The list of all files belonging to the 'CAS Bundle' is
%% given in the file `manifest.txt'.
%% 
%% Template article for cas-sc documentclass for 
%% single column output.

%\documentclass[a4paper,fleqn,longmktitle]{cas-sc}
\documentclass[a4paper,fleqn,colorlinks,hidelinks]{cas-sc}

%\usepackage[numbers]{natbib}
%\usepackage[authoryear]{natbib}
\usepackage[authoryear,longnamesfirst]{natbib}
\usepackage{graphicx}
\usepackage{float}
\usepackage{subcaption}
\usepackage{booktabs}
\usepackage{multirow}
% \usepackage[hidelinks]{hyperref}

%%%Author macros
% \def\tsc#1{\csdef{#1}{\textsc{\lowercase{#1}}\xspace}}
% \tsc{WGM}
% \tsc{QE}
% \tsc{EP}
% \tsc{PMS}
% \tsc{BEC}
% \tsc{DE}
%%%

\begin{document}
\let\WriteBookmarks\relax
\def\floatpagepagefraction{1}
\def\textpagefraction{.001}
\shorttitle{AtesN-DS: Offloading DNS into Hardware}
\shortauthors{Soares et~al.}
%\begin{frontmatter}

\title [mode = title]{AtesN-DS: Offloading DNS into Hardware with XDP}                      
% \tnotemark[1,2]

% \tnotetext[1]{This document is the results of the research
%    project funded by the National Science Foundation.}

% \tnotetext[2]{The second title footnote which is a longer text matter
%    to fill through the whole text width and overflow into
%    another line in the footnotes area of the first page.}


% ==========================================
% BLOCO DE AUTORES COM ORCID - UFMG
% ==========================================

% --- Autor 1 (Você) ---
\author[1]{João V. Soares S.}[orcid=0009-0002-7600-9784]
\ead{joaosoares@dcc.ufmg.br}
% \credit{Software, Data curation, Investigation, Writing - Original draft preparation}

% --- Autor 2 (André) ---
\author[1]{André G. Vieira}[orcid=0000-0002-7200-1825] % Substitua pelos números dele depois!
\ead{andregarcia@dcc.ufmg.br}
% \credit{Software, Methodology, Validation}

% --- Autor 3 (Luiz Felipe) ---
\author[1]{Luiz F. M. Vieira}[orcid=0000-0002-9050-3001]
\ead{lfvieira@dcc.ufmg.br}
% \credit{Conceptualization, Supervision, Writing - Review \& Editing}

% --- Autor 4 (Marcos Augusto - Correspondente) ---
\author[1]{Marcos A. M. Vieira}[orcid=0000-0001-9152-4114]
\cormark[1] 
\ead{mmvieira@dcc.ufmg.br}
% \credit{Conceptualization, Project administration, Supervision, Writing - Review \& Editing}

% --- Afiliação Única ---
\affiliation[1]{organization={Universidade Federal de Minas Gerais},
                department={Departamento de Ciência da Computação},
                addressline={Av. Antônio Carlos, 6627},
                city={Belo Horizonte},
                postcode={31270-901},
                state={MG},
                country={Brasil}}

% --- Texto do Rodapé do Autor Correspondente ---
\cortext[cor1]{Corresponding author}

% --- Nota de Agradecimento Geral ---
\nonumnote{This work was partially supported by CNPq.}

% ==========================================

\begin{abstract}
DNS is fundamental to the Internet, but traditional user-space resolvers incur high overhead from operating system network stack traversal. Kernel-space designs such as AtesN-DS reduce this cost by executing resolution in the eBPF/XDP driver layer, yet performance still degrades under heavy concurrency because of host interrupts and PCIe traffic. This paper presents a hybrid architecture that combines driver-space execution with programmable SmartNIC offloading. A hardware cache layer is instantiated in the NIC, where high-frequency cache hits are resolved at line rate without crossing the PCIe bus or triggering host interrupts, while cold misses cascade to the host resolution stack. To overcome hardware eBPF limitations—specifically the lack of native timers and LRU maps—we introduce a 1\,Hz host-synchronized epoch register and a host-DRAM Timing Wheel that executes $\mathcal{O}(1)$ Passive Evict-on-Write replacements with negligible host CPU consumption. Workload characterization on a 2-million-query production ISP trace demonstrates that our 56-byte hardware key layout accommodates the vast majority of real-world query traffic within on-chip SRAM. Evaluated on a physical 10GbE testbed under sustained open-loop traffic across four operational configurations, AtesN-DS sustains up to 3.53\,MQPS with sub-millisecond tail latency while reducing host CPU utilization by up to $29\times$ compared to driver-mode execution.
\end{abstract}

\begin{graphicalabstract}
\includegraphics{figs/cas-grabs.pdf}
\end{graphicalabstract}

\begin{highlights}
\item Hybrid DNS architecture combining XDP driver-space resolution with a SmartNIC hardware-offloaded SRAM cache.
\item Decoupled \textit{DNS Listener} and Passive Evict-on-Write timing wheel managing hardware tables with $0.053\%$ core overhead.
\item 64-byte cacheline-aligned hardware entry (56-byte key) capturing $97.90\%$ of production ISP query volume in wire format.
\item Evaluated across four operational modes (Unbound Standalone, Unbound + HW Cache, AtesN-DS Driver, and AtesN-DS Offloaded).
\item Delivers up to 3.53\,MQPS line-rate throughput and cuts host CPU utilization by $29\times$ under production ISP workloads.
\end{highlights}

\begin{keywords}
Domain Name System (DNS) \sep eBPF \sep XDP \sep Hardware offloading \sep Recursive resolver \sep Network stack bypass \sep Kernel-space processing
\end{keywords}


\maketitle


\section{Introduction}

The Domain Name System (DNS) is a distributed and hierarchical protocol originally specified in RFC~1034 and RFC~1035~\citep{mockapetris1987domain, rfc1035} that translates human-understandable domain names into machine-readable IP addresses. Currently, DNS is an indispensable part of the Internet; beyond being the pivot for name resolution, it also enables users to locate services~\citep{cheshire2013dns} and load balancing~\citep{brisco1995dns}. Furthermore, CDNs utilize DNS to route clients to the geographically closest servers based on their location~\citep{hours2016study} and to identify malicious botnets~\citep{choi2007botnet}. Thus, the global impact of DNS becomes evident, highlighting the importance of studies aimed at reducing name resolution time and the efficiency of this process.

The DNS infrastructure relies on recursive resolvers to close the gap between clients and authoritative servers. When a requested domain is missing from the local cache, the resolver must initiate an iterative traversal across the DNS hierarchy, a process that can trigger hundreds of nested messages for a single resolution in worst-case scenarios~\citep{afek2020nxnsattack}. Because traditional architectures implement these resolvers in user space~\citep{6932501}, every incoming and outgoing packet from this multi-hop sequence is forced to traverse the entire operating system network stack. This traversal triggers heavy performance penalties through continuous context switches and memory copies, consuming up to half of the available host CPU cycles even during baseline operations~\citep{10.1145/3672197.3673439}, which directly translates into tail-latency inflation for end-user applications~\citep{singla2014internet}.

To mitigate these bottlenecks, high-performance packet processing frameworks have increasingly explored kernel-bypass and programmable data-plane paradigms. Prominent user-space bypass frameworks, such as the Data Plane Development Kit (DPDK)~\citep{kourtis2015enhancing} and AF\_XDP sockets~\citep{10.1145/3281411.3281443}, accelerate packet ingestion by delivering frames directly to user space. However, these host-side frameworks present fundamental trade-offs: DPDK relies on aggressive Poll Mode Drivers (PMD) that monopolize 100\% of dedicated CPU cores in busy-waiting loops even when traffic is low, while completely severing integration with the Linux network stack and requiring basic networking protocols, routing, and firewalling to be laboriously re-implemented within application code. Furthermore, because both DPDK and AF\_XDP execute strictly on the host CPU, neither framework can synthesize an answer before the packet physically crosses the PCIe bus and consumes host memory bandwidth. In contrast, the extended Berkeley Packet Filter (eBPF) and the Express Data Path (XDP)~\citep{10.1145/3281411.3281443} provide an event-driven, kernel-cooperative execution model that preserves standard OS networking primitives while enabling packet handling directly at the driver layer without dedicated-core polling overhead~\citep{vieira2020fast, 9665095}. Crucially, XDP uniquely extends this paradigm to a hardware offload mode, enabling bytecode to execute directly inside the silicon of programmable SmartNICs.

Consequently, recent studies have shown that DNS resolvers placed directly in the XDP hook can achieve very high performance and substantially lower latency~\citep{sbesc_estendido, 10.1145/3672197.3673439}. However, leveraging XDP hardware offload for stateful DNS resolution introduces distinct architectural challenges. Because of strict SmartNIC silicon limits on instructions and memory, realizing an entire iterative resolver directly in the hardware data plane is impractical.

To overcome these constraints, this study presents a systems feasibility, architectural co-design, and empirical measurement study of DNS caching under the severe physical and verifier constraints of commodity SmartNIC eBPF hardware offload. We extend the AtesN-DS framework~\citep{sbesc_estendido} by combining its XDP driver-based resolver with a multi-tier cache structure built around XDP hardware offloading. By orchestrating a high-performance DNS cache directly on the programmable hardware alongside a secondary, larger cache layer in the driver space, the system intercepts and resolves the vast majority of queries at line rate in silicon while preserving seamless zero-hop fallback to host software.

In summary, the primary contributions of this work are threefold: $(i)$ a systems feasibility and architecture study of stateful DNS caching under the stringent constraints of commodity SmartNIC eBPF hardware offload, harmonizing in-silicon execution with kernel driver-space and user-space fallback; $(ii)$ an algorithmic and control-plane co-design introducing a 64-byte aligned key layout, a decoupled \textit{DNS Listener}, and a zero-PCIe-read $\mathcal{O}(1)$ Passive Evict-on-Write timing wheel that circumvents the lack of native hardware timers and LRU maps; and $(iii)$ a comprehensive empirical measurement study on a physical 10GbE testbed replaying a 2-million-query production ISP trace across four distinct operational modes, demonstrating up to 3.53\,MQPS line-rate throughput and a $29\times$ reduction in host CPU consumption.

The remainder of this paper is organized as follows. Section~\ref{sec:related_work} reviews the related work in the literature. Section~\ref{sec:proposed_system} details the hybrid architecture, multi-tier cache structure, and implementation details of the proposed system. Section~\ref{sec:evaluation} presents the experimental setup, workload characterization, and performance evaluation. Finally, Section~\ref{sec:conclusion} concludes the paper and discusses directions for future work.

\section{Literature Review}
\label{sec:related_work}

This section reviews core advancements in high-performance network services, kernel-embedded data planes, and in-network acceleration. We first describe the baseline AtesN-DS driver-mode recursive resolver in detail, establishing its operational mechanics and limitations. We then position our hybrid architecture against prior in-NIC and in-network acceleration systems, including FPGA-based responders, switch data-plane caches, in-kernel key-value caches, and pre-stack datagram offloaders. Finally, we contrast our contribution against the original conference paper (SBESC 2025) to delineate the new architectural extensions introduced in this journal version.

\subsection{The Baseline Driver-Mode AtesN-DS Recursive Resolver}
\label{sec:driver_mode_resolver}

The foundation of our work is the AtesN-DS recursive DNS resolver framework~\citep{sbesc_estendido}. Operating directly at the native driver-level eBPF/XDP hook (\texttt{XDP\_DRV}), AtesN-DS is designed to eliminate the severe operating system overheads inherent to traditional user-space resolvers. In conventional architectures, every incoming DNS query and outgoing authoritative resolution sequence must traverse the entire Linux network stack, incurring memory allocation penalties (\texttt{sk\_buff} allocation), SoftIRQ interrupt scheduling delays, and user-kernel context switching. Under high traffic volumes, these traversals consume substantial CPU cycles and induce tail-latency inflation.

To circumvent these bottlenecks, AtesN-DS embeds the resolution data path directly into the network device driver. The driver-mode resolver operates as follows:
\begin{enumerate}
    \item \textbf{Ingress Interception and Packet Parsing:} Upon packet arrival, the driver-space eBPF program intercepts UDP datagrams destined for port 53 directly from the network ring buffer. It validates Ethernet, IPv4, and UDP headers using direct memory pointers, verifying that the payload represents a well-formed DNS message. It parses the DNS header flags and extracts the Question section (query name, $\text{QTYPE}$, and $\text{QCLASS}$).
    \item \textbf{In-Driver Cache Hierarchy (\texttt{sw\_cache}):} The resolver queries an in-kernel BPF hash map (\texttt{sw\_cache}, preallocated with a capacity of up to 4,000,000 entries). Unlike hardware-constrained tables, \texttt{sw\_cache} resides in host main memory and supports full RFC~1035 domain names (up to 253 characters) and multiple resource records. Each entry stores the resolved IPv4/IPv6 addresses, record TTL, and an absolute expiration epoch timestamp. If a matching record is found and its TTL is valid, AtesN-DS immediately crafts a DNS response in-place, updates transport checksums, swaps source and destination MAC/IP addresses, and retransmits the frame directly back to the client interface using \texttt{XDP\_TX}, bypassing the upper host networking stack entirely.
    \item \textbf{Recursive Multi-Hop State Tracking (\texttt{queries\_map}):} On a cache miss, AtesN-DS must resolve the domain recursively without delegating to user space. To track in-flight asynchronous resolutions across multi-hop iterations (contacting root, TLD, and authoritative servers), the driver resolver allocates an in-flight transaction context within the \texttt{queries\_map} (sized for up to 4,587,576 concurrent queries). This table correlates the client's original transaction ID, IP address, and UDP port with a newly generated upstream transaction ID.
    \item \textbf{Upstream Zone Mapping (\texttt{ns\_cache}):} To minimize iterative traversal delays, AtesN-DS maintains an authoritative nameserver cache (\texttt{ns\_cache}, holding up to 8,000,000 zone records). This table maps domain zones and TLDs to known upstream nameserver IP addresses. If the target zone's nameserver is cached, AtesN-DS constructs an upstream iterative query and transmits it via \texttt{XDP\_TX} (or forwards it to the upstream WAN gateway). When the authoritative server returns a response, AtesN-DS intercepts the incoming packet, matches its transaction ID in \texttt{queries\_map}, updates \texttt{sw\_cache} with the newly learned A record, reconstitutes the client-facing response packet, and transmits it directly to the requesting host.
    \item \textbf{Proactive Cache Renewal (Prefetching):} To prevent latency spikes when popular domain records expire, AtesN-DS incorporates a proactive renewal mechanism. A lightweight user-space daemon monitors query access frequencies; when a frequently requested domain approaches its TTL expiration threshold, the daemon initiates a prefetch resolution to refresh the in-kernel cache before client requests experience a cold miss.
    \item \textbf{Selective Kernel and User-Space Fallback:} While AtesN-DS natively handles standard $\text{A}$/$\text{IN}$ iterative queries at line rate, handling complex or edge-case DNS protocol features (such as non-A query types, TCP connection fallbacks, DNSSEC signature validation, or EDNS0 large buffer negotiations) in eBPF would exceed kernel verifier instruction limits. Consequently, any packet exhibiting unhandled query types, malformed headers, or table saturation safely yields an \texttt{XDP\_PASS} action, allowing the frame to cascade cleanly into the standard host network stack and fall back to a full-featured user-space resolver (such as Unbound).
\end{enumerate}

While AtesN-DS in driver mode delivers substantial throughput gains over traditional user-space resolvers, its in-kernel stateful resolution is fundamentally bounded by host memory constraints and PCIe interrupt processing: every incoming packet must physically traverse the PCIe bus into host RAM, triggering NIC Receive Ring Buffer DMA transfers, PCIe MSI-X interrupts, and host SoftIRQ processing (\texttt{NET\_RX}). Specifically, tracking in-flight multi-hop recursive resolutions requires allocating transaction contexts within the in-kernel \texttt{recursive\_queries} map. When this table approaches memory exhaustion under extreme concurrency, AtesN-DS (consistent with hyDNS) implements a graceful degradation mechanism: queries that cannot be accommodated in driver state are passed down to the host networking stack via \texttt{XDP\_PASS}, falling back to a full-featured user-space resolver (such as Unbound) to avoid unhandled drops. Crucially, empirical evaluations reveal that cascading queries into user space under heavy load severely impairs performance rather than relieving host saturation. Transitioning frames to the host socket layer incurs substantial memory allocation penalties (\texttt{sk\_buff} allocation), intensifies SoftIRQ scheduling overhead, triggers user-kernel context switching, and induces severe contention in socket receive queues (\texttt{SO\_RCVBUF}), causing tail-latency inflation and socket-level packet drops. To quantify this degradation, our empirical stress evaluation of fallback ratios under an offered load of 500,000\,QPS ($n=30$) demonstrates that while pure in-driver processing ($0\%$ fallback) sustains $499{,}234$\,QPS with a 99th-percentile tail latency ($p_{99}$) of $4.19$\,ms, zero socket drops, and $79.56\%$ core utilization (dominated by SoftIRQ at $74.53\%$, with kernel system-call overhead at only $1.25\%$), routing even $25\%$ of queries to user space via \texttt{XDP\_PASS} immediately collapses throughput by $32.08\%$ (to $339{,}077$\,QPS). Simultaneously, host CPU Core~0 hits full saturation at $98.77\%$ (with kernel system overhead surging from $1.25\%$ to $28.78\%$ and user-space processing rising to $18.84\%$), tail latency ($p_{99}$) inflates by $5.04\times$ to $21.12$\,ms, and socket receive queues overflow, triggering over $110{,}000$ packet drops ($110{,}248 \pm 2{,}162$ \texttt{SO\_RCVBUF} drops). When all queries cascade to user space ($100\%$ fallback), throughput plummets by $76.54\%$ down to $117{,}127$\,QPS. Both the state capacity limits of in-driver recursive tracking and this steep empirical performance penalty of user-space fallback under load motivate offloading the primary line-rate caching tier directly into SmartNIC silicon.


\subsection{Positioning the Hardware Caching Layer Against Prior In-NIC and In-Network Systems}
\label{sec:positioning_prior_work}

To contextualize our hardware acceleration layer (\texttt{hw\_cache}), we directly contrast its in-silicon data-plane design against prominent prior hardware and pre-stack accelerators across programmable FPGAs, switch ASICs, and in-kernel caches. We emphasize that, consistent with prior in-network and in-NIC accelerators (such as Emu DNS, P4DNS, and NetCache), the SmartNIC hardware data plane of AtesN-DS operates strictly as a line-rate caching and answer-generation engine rather than an in-silicon recursive resolver. Maintaining recursive resolution state in hardware silicon is fundamentally constrained: iterative traversal requires tracking complex multi-hop dependency graphs, maintaining transaction contexts across hundreds of asynchronous sub-queries, reassembling out-of-order responses, and managing recursive timeout trees. Storing this state demands at least 256 bytes per active connection, far exceeding the SmartNIC firmware's 64-byte map entry ceiling, and requires thousands of stateful instructions that would exceed the NFP-4000's 1,536-instruction budget. By deliberately decoupling line-rate caching (handled entirely in SmartNIC silicon) from recursive resolution (delegated to the host driver and user space), AtesN-DS achieves an effective architectural balance: line-rate host bypass for high-frequency hits, with zero-network-hop fallback to full recursive capabilities on the host. Below, we compare our hardware caching tier against related data-plane systems:

\textbf{Emu DNS on NetFPGA:} \citet{203245} introduced \textit{Emu}, a framework for rapid prototyping of networking services on NetFPGA SUME boards using a C\#-to-FPGA high-level synthesis (HLS) compiler, implementing an authoritative DNS responder in FPGA silicon. While Emu demonstrates ultra-low packet latency, its standalone hardware design highlights key eviction and fallback trade-offs. In terms of cache eviction, Emu stores records in FPGA Block RAM (BRAM) without dynamic replacement; synthesizing fine-grained LRU or TTL-tracking state machines directly in FPGA logic is notoriously complex, consuming scarce logic slices (LUTs/DSPs) that would otherwise serve packet processing pipelines, thereby forcing reliance on static table allocations. Furthermore, implementing an automated fallback to host resolution software on an FPGA is technically cumbersome, requiring a custom PCIe DMA bridge and dedicated kernel drivers to be synthesized from scratch. In contrast, AtesN-DS avoids custom hardware synthesis by using standardized eBPF/XDP on a commodity SmartNIC (Netronome Agilio CX). Here, the tight SmartNIC-host co-design makes eviction policies simpler and more effective: rather than burning silicon logic on eviction hardware, AtesN-DS delegates eviction governance to a host DRAM circular timing wheel that selects victim slots in $\mathcal{O}(1)$ without PCIe reads via \textit{Passive Evict-on-Write}. Moreover, fallback is inherently native: any hardware miss returns \texttt{XDP\_PASS}, immediately handing the frame to the co-located host driver over standard DMA ring buffers with zero network hops.

\textbf{P4DNS and In-Network DNS Pipelines:} \citet{woodruff2019p4dns} proposed \textit{P4DNS}, mapping DNS query matching and response generation into programmable Match-Action Tables (MAT) targeted to NetFPGA-SUME and switch pipelines, while \citet{10619855} investigated application-layer header parsing in P4. While in-network FPGA and switch pipelines sustain line-rate forwarding, they face rigid state management, eviction, and fallback constraints. Feed-forward match-action pipelines cannot autonomously maintain dynamic linked lists or track second-resolution decaying TTLs in the data plane without controller intervention; eviction and record updates require an external control-plane agent or CPU to issue table-write RPCs, creating latency bottlenecks under dynamic DNS churn. Furthermore, because standalone FPGA or in-switch pipelines operate as separate network appliances, fallback on a cache miss cannot cascade over a host PCIe bus to a co-located resolver, forcing unhandled packets across additional network hops to backend servers. In contrast, AtesN-DS realizes an endhost SmartNIC-host co-design: misses cascade locally over PCIe via \texttt{XDP\_PASS} with zero network hops directly into the host driver, while eviction is governed asynchronously by the host daemon across the 32,768-entry SmartNIC cache, enabling deterministic on-demand replacement and minimum-TTL policies without burning FPGA logic slices or stalling the data plane.

\textbf{NetCache:} \citet{10.1145/3132747.3132764} developed \textit{NetCache}, an in-network key-value cache on Barefoot Tofino switch ASICs that absorbs load skews in distributed storage systems (e.g., Memcached) at over two billion queries per second. However, NetCache's hardware caching pipeline reveals fundamental differences in eviction and application scope compared to AtesN-DS. NetCache detects heavy hitters via data-plane sketches, but cannot evict items autonomously: cache replacement requires an external controller to aggregate statistics, compute victims, and push table updates down to the switch, introducing high reaction latency that is ill-suited for DNS where millions of entries expire with dynamic, second-resolution TTLs. Furthermore, NetCache is restricted to fixed 16-byte keys and 128-byte values capped at 64,000 entries in switch SRAM, unable to parse variable-length, multi-label DNS QNAMEs, and misses must be routed over the network to storage nodes. In contrast, AtesN-DS introduces an authentic multi-tier hierarchy: it leverages the server's SmartNIC adapter to host a compact, ultra-fast 32,768-entry Level-1 cache directly in on-chip SRAM fabric with up to 56-byte domain keys, while retaining a massive 4,000,000-entry Level-2 cache in host RAM. By synchronizing epoch time once per second via \texttt{time\_map}, the SmartNIC autonomously validates TTLs in silicon in $\mathcal{O}(1)$ with $0.04\,\mu\text{s}$ lookup latency, while table eviction is executed locally and on-demand by the host daemon upon insertion, eliminating off-path controller overhead entirely.

\textbf{BMC (Safe In-Kernel Caching):} \citet{performance:bmc} presented \textit{BMC}, an in-kernel cache for Memcached that uses eBPF at the driver-space XDP hook (\texttt{XDP\_DRV}) to serve UDP \texttt{GET} queries directly from kernel BPF maps, avoiding context switches and socket buffers. However, analyzing BMC illustrates why in-kernel software caching cannot match a hardware offload tier. Crucially, BMC relies on the Linux kernel's \texttt{BPF\_MAP\_TYPE\_LRU\_HASH} for automatic eviction; when porting eBPF to physical SmartNIC silicon (\texttt{XDP\_HW}), the hardware verifier and firmware strictly reject LRU maps due to the lack of hardware lock and linked-list primitives in silicon, preventing BMC's caching design from offloading to hardware as-is. Furthermore, because BMC executes strictly on host CPU cores, 100\% of incoming packets must cross the physical PCIe bus and generate host SoftIRQs (\texttt{NET\_RX}), consuming host CPU cycles under heavy load. AtesN-DS bridges this gap: by introducing \textit{Passive Evict-on-Write}, AtesN-DS enables standard hash maps (\texttt{BPF\_MAP\_TYPE\_HASH}) to function as self-governing caches in SmartNIC silicon without native LRU support, resolving over 99\% of queries directly in hardware before packets ever cross PCIe or interrupt the host CPU.

\textbf{In-Kernel Pre-Stack DNS Acceleration (Santa):} \citet{8493053} explicitly evaluated DNS offloading via \textit{Santa}, a pre-stack kernel rule processor that matches UDP datagrams early in the Linux kernel network stack, demonstrating a 5.5$\times$ throughput speedup for DNS over user-space BIND. However, like BMC, Santa is an in-kernel software module executing exclusively on host x86 CPU cores; every incoming DNS query must still traverse the physical PCIe bus and trigger host SoftIRQs (\texttt{NET\_RX}), incurring host memory bandwidth and interrupt overheads under high query volumes. Furthermore, Santa relies on static user-space rule tables without autonomous data-plane TTL decay validation or dynamic working-set replacement under memory pressure. AtesN-DS advances beyond pre-stack host software matching by moving the caching tier into physical SmartNIC silicon (\texttt{XDP\_HW}) and integrating an autonomous, time-synchronized eviction pipeline: by coupling an in-host DRAM circular timing wheel with atomic on-demand insertion over PCIe DMA, AtesN-DS answers hits in silicon with zero host CPU intervention while maintaining line-rate throughput and zero packet loss even under heavy working-set saturation.

\textbf{XDP-Accelerated Resolvers and Daemons (Knot DNS, NLnet Labs):} Recent production DNS projects have leveraged XDP and kernel bypass to enhance server throughput. CZ.NIC integrated native XDP support into Knot DNS~\citep{Vackova2021KnotXDP}, utilizing \texttt{AF\_XDP} zero-copy sockets and driver-mode XDP to deliver incoming UDP queries directly to user-space worker threads, achieving millions of queries per second for authoritative workloads. In parallel, NLnet Labs explored XDP for recursive and authoritative DNS in their XDPeriments project~\citep{toorop2021xdperiments}, developing driver-space eBPF modules for wire-speed packet filtering, rate limiting, and early cache answering. While these systems demonstrate the benefits of bypassing the OS network stack, they execute strictly as host-side software at the driver or user-space hook (\texttt{XDP\_DRV} or \texttt{AF\_XDP}). Consequently, 100\% of incoming queries traverse the physical PCIe bus, generate host CPU interrupts or monopolize cores in polling loops, and contend for host memory bandwidth. Neither Knot DNS nor NLnet Labs offloads DNS caching into programmable SmartNIC silicon (\texttt{XDP\_HW}), nor do they confront the stringent constraints of offloaded eBPF bytecode (such as the 1,536 native instruction budget, absence of dynamic memory allocation, and lack of hardware timers and LRU tables).

\textbf{In-Kernel DNS Acceleration (hyDNS):} \citet{10.1145/3672197.3673439} proposed \textit{hyDNS}, an in-kernel recursive DNS resolver combining an eBPF cache with recursive logic. In our preliminary conference work~\citep{sbesc_estendido}, we conducted a comparative evaluation against hyDNS, demonstrating that AtesN-DS significantly improves upon its design in functional capability, resilience, and wire-speed performance. Although hyDNS supports in-kernel recursion, it exhibits distinct functional deficiencies: (i)~it lacks a dynamically learned upstream zone delegation map (\texttt{ns\_cache}) naturally populated from referral responses; (ii)~it incorporates no proactive prefetching for expiring popular domain records to mitigate tail-latency spikes; (iii)~it provides no failure prevention or packet-loss resilience for UDP transport, whereas AtesN-DS dynamically detects zone lookup timeouts and dispatches queries across multiple diverse authoritative servers to mitigate UDP packet drops; and (iv)~it operates strictly as host software at the driver/TC layer, remaining bound to PCIe transfers and host CPU interrupts. Because our preliminary study established that AtesN-DS achieves higher throughput, lower tail latency, and greater operational resilience compared to hyDNS under identical test conditions, re-evaluating hyDNS on this SmartNIC testbed is redundant. Instead, to provide a rigorous, practical baseline for enterprise deployment, this work directly benchmarks the industry-standard user-space resolver Unbound alongside AtesN-DS.

\begin{table*}[t]
\centering
\small
\caption{Architectural taxonomy and technical comparison of AtesN-DS against prior in-network computing, in-NIC caching, and kernel-embedded accelerators.}
\label{tab:prior_work_comparison}
\begin{tabular}{@{}lllllll@{}}
\toprule
\textbf{System} & \textbf{Hardware Platform} & \textbf{Offload Tier} & \textbf{Key / Protocol Scope} & \textbf{Cache Eviction} & \textbf{Miss Fallback} & \textbf{Host CPU Cost} \\
\midrule
\textbf{Emu DNS} & NetFPGA SUME & FPGA Logic & Static Wire QNAME & Static (No Eviction) & Custom PCIe / None & None (Dedicated) \\
\textbf{P4DNS} & NetFPGA SUME (P4) & In-Network FPGA & Fixed-Field DNS & External Controller & Extra Network Hops & None (Off-Host) \\
\textbf{NetCache} & Barefoot Tofino & Switch MAT ASIC & Fixed 16\,B Key (Generic KV) & Off-Path Controller & Extra Network Hops & None (Off-Host) \\
\textbf{BMC} & Host x86 CPU & \texttt{XDP\_DRV} (Kernel) & Variable Key (Memcached) & Native Kernel LRU & In-Kernel Software & High (100\% PCIe/IRQ) \\
\textbf{Santa} & Host x86 CPU & Linux Kernel (Pre-Stack) & UDP Rule Tables (DNS) & Static Rule Updates & Linux Socket Layer & High (100\% PCIe/IRQ) \\
\textbf{Knot / NLnet} & Host x86 CPU & \texttt{XDP\_DRV} / \texttt{AF\_XDP} & RFC~1035 Wire Format & Kernel LRU / User-space & Host User Daemon & High (100\% PCIe/IRQ/Polling) \\
\textbf{hyDNS} & Host x86 CPU & \texttt{XDP\_DRV} / TC (Kernel) & RFC~1035 Wire Format & In-Kernel eBPF Map & In-Kernel / Sockets & High (100\% PCIe/IRQ) \\
\midrule
\textbf{AtesN-DS} & \textbf{Netronome SmartNIC} & \textbf{\texttt{XDP\_HW} (Silicon)} & \textbf{56\,B Wire QNAME} & \textbf{Passive Evict-on-Write} & \textbf{Co-Located PCIe} & \textbf{Near-Zero} \\
& \textbf{+ Host x86 CPU} & \textbf{+ \texttt{XDP\_DRV} (Kernel)} & \textbf{+ Dynamic Decaying TTL} & \textbf{+ Host Timing Wheel} & \textbf{(\texttt{XDP\_PASS}, 0 Hops)} & \textbf{(2.5\% at 1M QPS)} \\
\bottomrule
\end{tabular}
\end{table*}

\textbf{Synthesis of Core Novelty over Related Paradigms.} 
As codified in Table~\ref{tab:prior_work_comparison}, the high-level pattern of a data-plane cache answering hits and forwarding misses to host software has been explored in software caches (BMC) and in-network appliances (P4DNS). Rather than claiming this abstract pattern as novel, the core contribution of AtesN-DS lies in the \textbf{systems feasibility, architectural co-design, and empirical measurement study} of realizing stateful DNS caching under the concrete, severe physical constraints of commodity SmartNIC eBPF hardware offload. Specifically, AtesN-DS occupies a distinct design space defined by four architectural pillars:
\begin{enumerate}
    \item \textbf{Novelty over In-Kernel and Pre-Stack Accelerators (hyDNS, BMC, Santa, Knot DNS, NLnet Labs):} Prior software accelerators execute exclusively on host CPU cores. Although bypassing the socket layer via XDP or pre-stack hooks provides substantial speedups over user-space daemons, 100\% of incoming packets must still physically traverse the PCIe bus and trigger host SoftIRQ interrupts (\texttt{NET\_RX}) or monopolize cores in polling loops. As demonstrated in Section~\ref{sec:throughput_latency_eval}, sustaining 1\,MQPS in pure software driver mode consumes $72.6\%$ of host CPU cores. In contrast, AtesN-DS offloads the caching tier into physical SmartNIC silicon (\texttt{XDP\_HW}), resolving $>99\%$ of queries directly in hardware with zero PCIe transfers, zero host DMA, and zero CPU interrupts—reducing host CPU overhead by \textbf{29$\times$} (to only $2.5\%$). Furthermore, unlike hyDNS, AtesN-DS incorporates a dynamically learned authoritative zone map (\texttt{ns\_cache}), proactive cache prefetching, and multi-server UDP retry fanout for loss resilience.
    \item \textbf{Novelty over In-Network Switch and FPGA Pipelines (P4DNS, NetCache, Emu DNS):} Dedicated switch and FPGA accelerators reside as off-host devices or custom hardware boards within the network fabric. Consequently, cache misses cannot be resolved locally and must traverse additional physical network hops to reach backend application servers, compounding tail latency under cache churn. Moreover, feed-forward match-action switch and FPGA pipelines cannot autonomously execute timer expiration or linked-list manipulations in data-plane silicon, forcing reliance on slow, out-of-band controller RPCs. AtesN-DS leverages an endhost SmartNIC co-design: hardware misses cascade locally over PCIe via \texttt{XDP\_PASS} with zero network hops directly into the co-located recursive engine, while cache consistency and eviction are maintained locally and asynchronously by the host runtime.
    \item \textbf{Circumventing Fundamental eBPF Hardware Offload Limitations:} Unlike software eBPF frameworks like BMC that rely on \texttt{BPF\_MAP\_TYPE\_LRU\_HASH} (which SmartNIC firmware and verifiers strictly reject), AtesN-DS demonstrates how to operate standard hash maps as self-governing caches in silicon. By introducing the 1\,Hz host-synchronized epoch register (\texttt{time\_map}) and the zero-PCIe-read \textit{Passive Evict-on-Write} timing wheel, AtesN-DS solves the dual hardware barriers of missing silicon timers and absent LRU maps with negligible host CPU overhead ($0.053\%$).
    \item \textbf{Domain-Specific Wire Safety and Trace-Driven Evaluation:} Generic key-value offloads (NetCache, BMC) operate on fixed-size integer keys or opaque hashes with static values, serving uniform \texttt{GET} primitives without protocol-specific semantics. In contrast, DNS presents variable-length hierarchical label encoding, dynamic per-record decaying TTLs, Transaction ID matching, in-place header flag synthesizing, name compression pointers ($0xc00c$), and vulnerability to cache-poisoning attacks. AtesN-DS implements full DNS wire-format parsing directly in bytecode, synchronizes real time at 1\,Hz, isolates cache ingestion behind the decoupled, stateful \textit{DNS Listener} to enforce transaction safety, and provides the first exhaustive physical testbed measurement study replaying a production ISP trace under varying offered rates.
\end{enumerate}

\subsection{Novelty and Extensions Relative to the SBESC 2025 Paper}
\label{sec:novelty_sbesc}

The preliminary version of AtesN-DS presented at SBESC 2025~\citep{sbesc_estendido} introduced the concept of an in-kernel recursive DNS resolver executing entirely in software at the native XDP driver layer (\texttt{XDP\_DRV}). While the conference paper established the viability of kernel-space recursive resolution, it operated exclusively on the host CPU, meaning that all traffic traversed the physical PCIe bus and generated host SoftIRQs. 

This journal article significantly expands upon the SBESC 2025 conference paper through five major architectural, algorithmic, and empirical contributions:
\begin{enumerate}
    \item \textbf{Hardware Data-Plane Offloading (\texttt{XDP\_HW}):} We design and implement the first hybrid DNS architecture that offloads stateful caching directly into the silicon of a programmable SmartNIC (Netronome Agilio CX 2$\times$10GbE), enabling line-rate cache hits to be answered in hardware with zero PCIe transfers, zero host DMA, and zero host CPU interrupts.
    \item \textbf{Optimized 56-Byte Key and 64-Byte Aligned Cache Architecture:} To overcome the strict physical memory and instruction limits of network microengines, we restructured the offloaded packet inspection pipeline into an optimized, two-phase bounded unrolling scheme. This reduced the compiled eBPF bytecode to \textbf{589 native instructions} (utilizing just 38.3\% of the 1,536-instruction budget) and enabled an optimal 64-byte entry layout (56-byte domain key + 8-byte value) that accommodates the 32,768 working-set records directly within $\approx 2.35$\,MB of ultra-fast on-chip SmartNIC SRAM fabric with natural memory alignment and sub-microsecond in-silicon lookups.
    \item \textbf{Zero-PCIe-Read Passive Eviction and Host Timing Wheel:} Because SmartNIC silicon lacks native LRU map support (\texttt{BPF\_MAP\_TYPE\_LRU\_HASH}), naively polling hardware tables over PCIe to evict expired records triggers severe interconnect saturation and hash bucket locking. We introduce the \textit{Passive Evict-on-Write Daemon} and a host-side 86,400-bucket circular timing wheel that tracks domain expirations with zero PCIe read traffic, executing $\mathcal{O}(1)$ on-demand replacements with negligible host CPU overhead ($0.053\%$).
    \item \textbf{Asymmetric, Tamper-Proof Cache Population (DNS Listener):} To prevent unauthenticated or spoofed packets on the wire from directly altering hardware state, we introduce the decoupled \textit{DNS Listener}. Operating as a verified eBPF filter at the host driver hook, it statefully matches authoritative responses against in-flight queries before feeding valid A records into an asynchronous host ring buffer, which a privileged daemon commits over PCIe DMA. This design ensures complete cache-poisoning protection and allows the hardware cache to be deployed transparently in front of standard upstream resolvers (e.g., BIND or Unbound) without source code modification.
    \item \textbf{Production ISP Trace Evaluation and Tail-Latency Analysis:} While the preliminary conference paper evaluated synthetic uniform traffic, this journal work conducts exhaustive empirical evaluations on a physical 10GbE point-to-point testbed replaying a large-scale production ISP DNS resolver trace of \citet{11314363}. We expand the evaluation scope to report tail latency percentiles ($p_{50}$, $p_{90}$, $p_{99}$, $p_{99.9}$), packet loss, host SoftIRQ reduction, per-core CPU saturation, and SmartNIC hardware resource utilization across load rates up to 5,000,000 queries per second.
\end{enumerate}

\subsection{Latency, Caching, and eBPF in Modern Systems}
\label{sec:background_latency_caching}

Several studies have investigated the critical impact of network latency on user experience~\citep{attig2017system}, including extensive research on web browsing~\citep{arapakis2021impact} and cloud-based applications~\citep{claypool2014effects}. Among these works, some identify the operating system network stack as one of the primary bottlenecks responsible for latency inflation~\citep{singla2014internet}. To overcome this architectural limitation, researchers have developed highly efficient network stacks~\citep{10.1145/2740070.2626311} as well as stack-bypass strategies~\citep{10.1145/3297156.3297242}. Notable paradigms among these bypass techniques include the Data Plane Development Kit (DPDK)~\citep{kourtis2015enhancing} and the adoption of eBPF/XDP~\citep{cai2021understanding}.

Beyond bypassing network stack traversals, caching represents another well-established strategy for latency reduction, with eBPF emerging as a highly effective enabler for in-kernel cache implementations. Frameworks such as \texttt{cache\_ext}~\citep{zussman2025cache_ext} and LearnedCache~\citep{qi2026learnedcache} (which leverages perceptron-based eviction policies) demonstrate the effectiveness and versatility of eBPF in customizing OS page cache management, while MiddleCache~\citep{pang2023middlecache} extends in-kernel caching to support TCP-based queries.

Beyond packet filtering, recent works have leveraged eBPF to bypass upper-layer kernel overheads in both storage and transport subsystems. For instance, XRP~\citep{280870} enables user-defined storage functions, such as index lookups and aggregations to execute directly within an eBPF hook in the NVMe driver, drastically reducing OS storage stack overheads for microsecond-scale devices. Similarly, eTran~\citep{chen2025etran} extends eBPF to make the kernel networking stack agilely extensible, allowing developers to safely implement and customize complex transport designs, such as TCP congestion controls and clean-slate protocols like Homa, directly within the Linux kernel while achieving near user-space performance. As eBPF adoption expands across diverse system domains, ensuring execution isolation and multi-tenant security becomes critical, addressed by vBPF~\citep{zhang2026virtualizing} via late-binding virtualization and KrakenGuard~\citep{patel2026krakenguard} via symbolic execution at load time. Furthermore, \citet{9665095} demonstrate that leveraging eBPF with XDP for network monitoring yields higher computational efficiency compared to kernel-bypass alternatives like DPDK. Moving beyond general-purpose host CPUs, platforms such as hXDP~\citep{performance:hxdp} and eBPFlow~\citep{performance:ebpflow,networking:application_layer} demonstrate that offloading packet processing to FPGAs programmed with eBPF can achieve ultra-low latency and high throughput.

\section{Proposed System}
\label{sec:proposed_system}

This section describes the proposed multi-tier DNS design and implementation, which partitions packet processing across programmable SmartNIC silicon, the kernel driver, and user space. Rather than viewing data-plane caching as an abstract pattern, this work addresses the concrete systems challenge of realizing stateful DNS caching under the severe physical and verifier constraints of commodity eBPF hardware offload (\texttt{XDP\_HW}). To achieve line-rate acceleration within these hardware limits, we co-design the in-silicon data plane with host-level control mechanisms: resolving the absence of hardware timers via a 1\,Hz epoch register, mitigating the lack of native hardware LRU maps via a zero-PCIe-read \textit{Passive Evict-on-Write} timing wheel, and enforcing cache integrity via an asymmetric, decoupled \textit{DNS Listener}.

\subsection{Design}

To maximize packet processing throughput while working around the resource constraints of network hardware, our system adopts a hybrid, multi-tier resolution architecture (depicted in Figure~\ref{fig:1}). Instead of processing all DNS logic within a single execution tier, the system decouples high-frequency caching from full recursive resolution across three specialized operational levels: an offloaded hardware data plane in the SmartNIC silicon, an in-kernel driver resolver, and a user-space control and fallback tier. 

The core of the hardware offload layer operates as a line-rate, read-mostly cache executing directly within the NIC silicon, dedicated to accelerating client query resolution while decoupling cache population to authenticated host control channels. For incoming client DNS queries, the hardware cache inspects the requested domain name. On a cache hit with a valid TTL, the NIC directly constructs the DNS answer, rewrites transport headers, and retransmits the frame back to the client at line rate (\texttt{XDP\_TX}), completely bypassing host interrupts, PCIe bus traversals, and upper-stack overhead. Conversely, on a cache miss or TTL expiration, the query cleanly cascades down over PCIe via \texttt{XDP\_PASS} to the native AtesN-DS driver layer for high-throughput recursive resolution or to user space under severe traffic overloads. Crucially, the hardware data plane enforces strict immutability against wire packets: it contains no write or insertion logic for external responses. Incoming DNS responses arriving from the network wire immediately yield \texttt{XDP\_PASS} to the host stack, ensuring that raw network frames cannot directly alter hardware memory. As detailed below, cache population is handled exclusively through an asymmetric control-plane pipeline fed by verified host resolutions.

To maintain hardware simplicity and keep the offloaded data plane lightweight, the architecture delegates state maintenance and cache lifecycle functions to dedicated user-space control daemons running on the host. First, a \textit{time daemon} periodically synchronizes the host epoch clock with the hardware memory, enabling the offloaded data plane to perform absolute TTL freshness validation during packet processing without requiring native hardware timer support. Second, to manage cache saturation and stale record reclamation in the absence of native hardware LRU support, we developed the \textit{Passive Evict-on-Write Daemon}, which couples an in-host DRAM timing wheel with reactive on-demand eviction directly tied to record insertion. As detailed in Section~\ref{sec:implementation} and empirically demonstrated in Section~\ref{sec:saturation_eval}, this passive architecture sustains 100\% hardware cache capacity without artificial headroom buffers, operates with negligible host CPU overhead, and completely eliminates PCIe interconnect storms and hash bucket spinlock contention on the SmartNIC datapath.

As shown in Figure~\ref{fig:1}, dark blue highlights the components of the hardware-offloaded cache subsystem, encompassing both the in-NIC data plane and the host-side control daemons, and gray represents the native driver-space AtesN-DS resolver stack.

The proposed hardware cache architecture cleanly decouples high-speed query answering from recursive resolution. To enable modular and secure cache population, we introduce the \textit{DNS Listener}, a lightweight eBPF program attached to the native (driver) or generic (emulated) XDP hook. The DNS Listener operates in two deployment modes: (1)~\textit{Integrated Mode}, embedded directly within the in-kernel AtesN-DS resolver, or (2)~\textit{Standalone Mode}, functioning as an autonomous network filter positioned transparently in front of any standard upstream resolver (such as BIND, Unbound, or PowerDNS). By intercepting outgoing queries and incoming authoritative responses at the XDP layer, the DNS Listener statefully tracks in-flight DNS transactions (recording transaction IDs, randomized UDP ports, and domain names). Only when an incoming response strictly matches an outstanding in-flight query does the DNS Listener extract the validated A record and dispatch an event to the host ring buffer (\texttt{hw\_sync\_ringbuf}). A privileged host user-space daemon subsequently commits the record to the SmartNIC lookup table via \texttt{bpf\_map\_update\_elem} over PCIe DMA. This design delivers strong architectural security and complete modularity: unauthenticated or spoofed packets from the network wire cannot bypass resolution validation or pollute the in-silicon lookup table, and the hardware caching layer can be seamlessly coupled with any existing DNS server software without modifying its internal source code.

Because these operational tiers execute in complete isolation without sharing state or parsed packet metadata, each layer must independently parse the DNS query header upon packet arrival. In the event of a cache miss at lower hardware layers, this decoupled inspection introduces a negligible parsing overhead, as upper host layers must re-parse the domain upon receiving the forwarded packet. However, this trade-off yields substantial performance dividends: on a cache hit, the hardware data plane resolves the query entirely in silicon, completely bypassing host interrupts, PCIe bus traversal, and upper-stack overhead. Ultimately, this isolated design was chosen for architectural simplicity, ensuring modularity, robust fallback mechanisms, and zero modification requirements for existing upstream resolution software.

\begin{figure}
	\centering
	\includegraphics[width=.9\textwidth]{figs/atesnds.drawio-hw-7.drawio.png}
	\caption{AtesN-DS Architecture with Offloaded Cache. (1) Indicates packets or unhandled DNS queries forwarded to user space due to architectural limitations or hardware saturation. (2) Represents packets generated in user space for proactive cache renewal of popular domains or fault prevention. (3) Denotes queries fully resolved by the traditional user-space resolver.}
	\label{fig:1}
\end{figure}


\subsection{Implementation}
\label{sec:implementation}

The hardware-level caching layer was implemented using the eBPF/XDP offload framework, attaching the data plane execution directly to the SmartNIC silicon. While eBPF/XDP offloading has emerged as a compelling paradigm for line-rate packet processing, adapting high-level stateful application logic such as a DNS lookup table to network interface hardware introduces two fundamental architectural challenges: time abstraction and restricted map operations. The overall architecture and execution pipeline of this hardware-level caching mechanism are illustrated in Figure~\ref{fig:hw_pipeline}.

First, because the XDP offload engine lacks time-based helpers such as \texttt{bpf\_ktime\_get\_ns()}, the hardware data plane has no intrinsic notion of time to evaluate domain TTL expiration. To bridge the host clock to the SmartNIC, the architecture employs a single-element \texttt{BPF\_MAP\_TYPE\_ARRAY} named \texttt{time\_map}. Standard eBPF array maps mandate a 32-bit unsigned integer as the key index. The map value is also a 4-byte unsigned integer holding the global epoch timestamp in seconds. A lightweight host-side \textit{time daemon} (Figure~\ref{fig:1}) periodically writes the current host epoch timestamp to index zero once per second ($1\,\text{Hz}$), providing fine-grained 1-second TTL accuracy with negligible host overhead. During packet lookup, the offloaded program performs an $\mathcal{O}(1)$ query against index zero, enabling in-NIC TTL validation by comparing the retrieved epoch against stored record expiration timestamps.

Second, SmartNIC hardware architectures inherently lack silicon support for native Least Recently Used (LRU) maps (\texttt{BPF\_MAP\_TYPE\_LRU\_HASH}), which would otherwise manage record eviction automatically upon table saturation. This hardware limitation forces the caching mechanism to rely on a standard hash map (\texttt{BPF\_MAP\_TYPE\_HASH}), instantiated within the offloaded program as \texttt{hw\_cache} (designated \texttt{level\_one\_cache} in firmware). To achieve sub-microsecond in-silicon lookups without crossing external memory controllers, \texttt{hw\_cache} is explicitly provisioned with $32,768$ entries ($2^{15}$), fitting strictly within the SmartNIC's ultra-low latency on-chip SRAM fabric (internal SRAM and Cluster Local Scratchpad, $\approx 3.5$\,MB total capacity). The map entry is structured with a 56-byte key, accommodating fully qualified domain names up to 56 characters, and an 8-byte value containing two 32-bit unsigned integers: a 4-byte IPv4 response address and a 4-byte absolute expiration timestamp. Crucially, this 56-byte key combined with the 8-byte value yields a total entry size of exactly 64 bytes, perfectly matching the hardware's 64-byte memory alignment boundary and consuming $\approx 2.35$\,MB of on-chip SRAM including hash bucket pointers.

Empirical evaluation on our production ISP resolver trace ($2{,}000{,}000$ queries across $83{,}659$ unique domains) confirms that $97.90\%$ of all incoming queries ($76.20\%$ of unique domains) have RFC~1035 wire-format domain names $\le 56$ bytes (with a median wire length of $21.0$\,B and a mean of $24.81$\,B across replayed traffic). Consequently, this 56-byte key layout achieves an optimal design point: it seamlessly accommodates practically all real-world DNS traffic directly within on-chip SRAM, resolving heavy-hitters in $0.04\,\mu\text{s}$ while delegating only the $2.10\%$ long-name tail to the host cache.

Because standard \texttt{BPF\_MAP\_TYPE\_HASH} tables cannot autonomously purge stale entries when capacity is reached, managing the lifecycle of cached records relies on a coordinated dual-plane operational strategy. At the hardware data plane, the offloaded XDP program performs inline lazy TTL evaluation, checking entry freshness during packet lookup and triggering immediate $\mathcal{O}(1)$ deletions via \texttt{bpf\_map\_delete\_elem} upon encountering expired records. However, because low-frequency or cold stale records may not be accessed again by client queries, relying solely on data-plane lookups would eventually lead to capacity exhaustion. 

To address hardware capacity limits and ensure that no newly resolved domain records are lost upon reaching table saturation, the system integrates the Passive Evict-on-Write Daemon. Instead of performing costly periodic sweeps across the PCIe bus, the host control plane maintains an in-DRAM \textit{Timing Wheel}, a hierarchical, slotted circular array that tracks record expiration lifetimes with zero PCIe interaction.

In this passive paradigm, background sweeps across the network card are completely eliminated. Eviction is triggered strictly on demand upon inserting a newly resolved DNS record. During an unsaturated state, when the hardware table has not reached capacity, the ingestion worker immediately issues an $\mathcal{O}(1)$ \texttt{bpf\_map\_update\_elem} over PCIe DMA and registers the key alongside its expiration epoch in the host DRAM timing wheel. Conversely, during a saturated state, when the hardware table is full, the ingestion worker executes a deterministic $\mathcal{O}(1)$ 1-for-1 demand replacement: it retrieves a victim from the Timing Wheel, which prioritizes popping already expired records or, if none are yet expired, the record with the minimum remaining TTL. The worker then issues an individual \texttt{bpf\_map\_delete\_elem} to evict the victim from the SmartNIC, immediately followed by \texttt{bpf\_map\_update\_elem} to admit the new entry.

From an execution perspective, the packet processing pipeline branches based on DNS message headers and protocol semantics evaluated inline at the SmartNIC ingress. Upon packet arrival, the offloaded XDP program verifies that the frame represents an IPv4 UDP packet directed to port 53. To enforce strict type isolation and prevent cache pollution across different resource classes, the parser explicitly validates that the query contains a single question with $\text{QTYPE}=\text{A}$ and $\text{QCLASS}=\text{IN}$. Any query requesting alternative record types (such as \texttt{AAAA}, \texttt{MX}, \texttt{TXT}, \texttt{SRV}, or \texttt{CNAME}), non-IN classes, or advertising large EDNS0 buffers immediately returns \texttt{XDP\_PASS}, cascading untouched to the host resolver stack. Also, queries requiring full multi-IP RRsets or DNSSEC cryptographic validation ($\text{DO}=1$) are handled in host software. By verifying $\text{A}/\text{IN}$ as a mandatory precondition in bytecode, the hardware lookup table can dedicate its 56-byte key budget exclusively to the domain name, avoiding wasting key bytes on static type and class identifiers within the SmartNIC's hardware memory layout.

For eligible A-record queries, the program retrieves the current epoch from \texttt{time\_map} and executes an $\mathcal{O}(1)$ lookup against \texttt{hw\_cache} using the extracted \texttt{QNAME}. On a cache hit, the hardware validates freshness by dynamically calculating the decaying remaining TTL. If the remaining TTL is less than $1\,\text{s}$, the record is deemed expired, triggering inline eviction via \texttt{bpf\_map\_delete\_elem} and forwarding the query to the host via \texttt{XDP\_PASS}. Otherwise, the hardware engine synthesizes the DNS answer entirely in silicon: it preserves the client's 16-bit Transaction ID in place, sets response flags to $\text{QR}=1$, $\text{RD}=1$, $\text{RA}=1$, $\text{RCODE}=\text{NOERROR}$, updates transport headers with incrementally recalculated checksums, and appends an answer record populated with the remaining TTL and the resolved 4-byte IPv4 address. Crucially, the Question section is left byte-for-byte intact in packet memory, and the answer record utilizes standard 2-byte DNS name compression ($0xc00c$) pointing directly to the original Question name at offset 12. The program then returns \texttt{XDP\_TX} to retransmit the synthesized response immediately at line rate, completely bypassing host interrupts and PCIe traversal.

Conversely, on a cache miss, or when processing incoming authoritative responses ($\text{QR}=1$) from upstream nameservers, the SmartNIC executes strictly zero write operations and immediately yields \texttt{XDP\_PASS} to hand the packet to the host. Because the XDP offload hook operates on the network interface's RX (ingress) path, the hardware cache is never populated from host-to-client TX transmissions. Instead, cache ingestion is governed by the \textit{DNS Listener} eBPF program, attached to the host's native or generic XDP layer. As recursive queries are dispatched upstream, the DNS Listener records active in-flight transactions in an eBPF tracking map (\texttt{recursive\_queries}), tracking the 16-bit transaction ID ($\text{TXID}$), the randomized ephemeral UDP destination port, and the requested domain name ($\text{QNAME}$). When an authoritative response arrives on RX, the DNS Listener validates it against the active in-flight state. Only upon verifying an exact match across all three tuples does it extract the validated record and dispatch an event to the \texttt{hw\_sync\_ringbuf} ring buffer. The user-space ingestion daemon (\texttt{daemon\_ingest.c}) then consumes the event, updates the host DRAM timing wheel, and issues an authorized $\mathcal{O}(1)$ \texttt{bpf\_map\_update\_elem} via PCIe DMA to commit the entry into the SmartNIC. This decoupled architecture ensures that external network packets can never directly mutate the SmartNIC cache, providing robust protocol conformance and end-to-end security whether paired with AtesN-DS or a third-party recursive resolver.

\begin{figure}[h]
    \centering
    \includegraphics[height=0.5\textheight]{figs/ChatGPT Image Aug 6, 2026, 11_05_05 PM.png}
    \caption{Hardware-level eBPF/XDP packet processing pipeline within the SmartNIC data plane. Incoming DNS queries undergo domain extraction and time validation via \texttt{time\_map}. Valid cache hits are resolved entirely in silicon and retransmitted immediately via \texttt{XDP\_TX} (green path). On a cache miss, complex query, or TTL expiration, the frame is forwarded to the host stack via \texttt{XDP\_PASS}. Incoming DNS responses on the wire bypass the NIC via \texttt{XDP\_PASS} for stateful in-flight verification by the DNS Listener (integrated or standalone); verified records are subsequently ingested into \texttt{hw\_cache} by the host control daemon via PCIe DMA (orange path).}
    \label{fig:hw_pipeline}
\end{figure}

\subsection{Architectural Trade-offs: Hardware Offloading and In-NIC Scope}
\label{sec:tradeoffs}

When designing high-throughput DNS architectures, user-space bypass frameworks such as DPDK~\citep{kourtis2015enhancing} and AF\_XDP sockets~\citep{10.1145/3281411.3281443} represent well-known alternatives. However, selecting XDP hardware offloading over these host-side paradigms was guided by three key architectural considerations. First, the primary research objective of this work is to evaluate packet processing executed directly within programmable SmartNIC silicon, which neither DPDK nor AF\_XDP permits. Crucially, no host-side framework can answer a query before it physically crosses the PCIe bus and consumes host memory bandwidth. Second, DPDK relies on aggressive Poll Mode Drivers (PMD) that monopolize 100\% of dedicated CPU cores in busy-waiting loops even under low traffic, contradicting our goal of liberating host resources for recursive resolution. Third, DPDK's complete kernel bypass severs integration with the OS network stack, requiring routing, firewalling (iptables/nftables), and protocol edge cases (such as TCP fallbacks for truncated responses) to be entirely re-implemented in application code. In contrast, eBPF/XDP hardware offloading resolves cache hits in silicon with zero host CPU datapath overhead, while preserving a cooperative model where cache misses, non-A queries, and long-tail packets fall through cleanly via \texttt{XDP\_PASS} to the native Linux network stack and existing resolver software. While running eBPF firmware on the SmartNIC can supersede certain fixed-function ASIC pipelines (such as generic L4 checksumming), absorbing millions of queries in silicon with sub-microsecond latency provides far greater system relief.

A complementary architectural question is \textit{why offload only a caching layer instead of a fully-fledged recursive resolver into hardware?} From a workload perspective, real-world DNS traffic follows a highly skewed, Zipfian distribution heavily concentrated around a small subset of high-frequency domains~\citep{10.1145/505202.505223}. Because the vast majority of incoming queries target these popular records, offloading a lightweight caching layer alone captures nearly all potential query hits where line-rate, sub-microsecond acceleration provides systemic relief to the host. Conversely, attempting to push full recursive resolution into silicon yields diminishing returns. Recursive resolution is an inherently complex, long-running operation involving pending socket tracking, intermediate referral parsing, and timeout retries across authoritative servers. Accelerating these long-tail cache misses in hardware provides negligible performance gains, as their latency is overwhelmingly dominated by wide-area network (WAN) round-trip times (RTTs) rather than host processing overhead.

Furthermore, attempting to implement stateful recursive logic is fundamentally obstructed by the low-level constraints of eBPF/XDP offloading established in Section~\ref{sec:implementation}. As highlighted in prior eBPF offload studies~\citep{tibav:41657} and confirmed by our implementation, SmartNIC execution environments enforce strict silicon and verifier boundaries. The 4,096 assembly-like instruction budget and the lack of program chaining via tail calls (\texttt{bpf\_tail\_call}) prevent decomposing complex recursive state machines into modular BPF programs. Most critically, maintaining active tracking contexts for recursive queries demands at least 256 bytes per record, fundamentally exceeding the entry footprint and memory constraints enforced by SmartNIC eBPF map offloading. 

Consequently, our implementation proves that pairing a high-throughput hardware cache for line-rate hit absorption with host-based software recursive fallback represents the optimal architectural trade-off between hardware performance and software flexibility.

\subsection{Realistic DNS Cases, Cache Poisoning Protection, and Operational Scope}
\label{sec:threat_model}

A critical architectural consideration for in-silicon DNS acceleration is how the system handles diverse protocol semantics, attack vectors, and operational edge cases without sacrificing hardware performance or protocol conformance. To reconcile line-rate acceleration with the vast complexity of RFC standards, our architecture establishes a strict division of responsibility: the SmartNIC data plane is specialized as a line-rate fast-path cache strictly for $\text{QTYPE}=\text{A}$ and $\text{QCLASS}=\text{IN}$ records up to 56 bytes, while all complex, non-A, or long-tail DNS behaviors gracefully cascade to the host driver and user-space resolver stack via non-disruptive \texttt{XDP\_PASS} fallback.

\textbf{Specialization to IPv4 A/IN Records and Single-Record Answers.} Specializing the hardware table exclusively to $\text{A}/\text{IN}$ queries represents a deliberate design decision guided by empirical traffic distributions: IPv4 A-record queries constitute over $85\%$ to $90\%$ of enterprise volume~\cite{SINGH201914}. By evaluating $\text{QTYPE}=1$ and $\text{QCLASS}=1$ inline at the packet parser, the lookup table dispenses with multi-byte type and class fields, dedicating its key budget entirely to the domain name. This enables the compact 64-byte entry layout ($56\,\text{B}\text{ key} + 8\,\text{B}\text{ value}$) matching the hardware's 64-byte alignment boundary and fitting within the 1,536 native NFP instruction limit. Crucially, on a hardware cache hit, the SmartNIC synthesizes and returns \textbf{exactly one A record} (a single 4-byte IPv4 address) populated with the decaying remaining TTL and utilizing 2-byte DNS compression ($0xc00c$) pointing to the Question section. If an upstream domain delegates a multi-IP RRset (e.g., round-robin pools), the hardware cache indexes the primary validated A record, whereas multi-record expansion or alternative classes cascade to the host stack.

\textbf{Graceful Fallback for Complex DNS Query Types.} Rather than dropping or rejecting unsupported query formats, the offloaded ingress parser inspects the DNS question headers and immediately returns \texttt{XDP\_PASS} with zero packet loss for the following cases:
\begin{itemize}
    \item \textbf{Long Domain Names ($>56$ bytes):} Queries targeting domain names exceeding 56 characters cannot be indexed in the 56-byte hardware table. These queries bypass hardware lookup and cascade untouched to the host driver-space resolver.
    \item \textbf{CNAME Aliases and Alternative Query Types (\texttt{AAAA}, \texttt{MX}, \texttt{TXT}, \texttt{SRV}, \texttt{PTR}):} Incoming queries requesting IPv6 addresses (\texttt{AAAA}), mail routing (\texttt{MX}), or text records (\texttt{TXT}) immediately pass to host software via \texttt{XDP\_PASS}. Specifically, resolving Canonical Name (\texttt{CNAME}) aliases requires following recursive redirection chains across multiple zones, which fundamentally exceeds the microengine's single-lookup feed-forward instruction budget; consequently, CNAME traversal is handled entirely in host software.
    \item \textbf{DNSSEC Cryptographic Validation ($\text{DO}=1$):} Cryptographic signature verification (\texttt{RRSIG}, \texttt{DNSKEY}, \texttt{DS}) requires complex asymmetric cryptographic primitives, dynamic chain traversal, and large payload handling incompatible with SmartNIC microcode budgets. All queries advertising the DNSSEC OK ($\text{DO}$) flag bypass the hardware cache to be validated securely in host software.
    \item \textbf{EDNS0 Options and Large Payloads ($>512$ bytes):} When a query includes an EDNS0 \texttt{OPT} pseudo-RR~\citep[RFC~6891]{rfc6891}, advertises custom option codes (such as EDNS Client Subnet, ECS, or DNS Cookies), or requests buffer sizes $>512$ bytes, the parser immediately yields \texttt{XDP\_PASS}. This ensures that EDNS0 negotiation and option processing are executed with complete RFC conformance by host software, while preserving a strictly bounded $\approx 1.0\times$ response-to-query ratio in silicon to prevent DNS amplification attacks.
    \item \textbf{Stateful TCP Fallback:} The SmartNIC offload hook strictly filters stateless UDP frames on port 53. If a response exceeds the maximum transmission unit or triggers a truncation flag ($\text{TC}=1$), the client reissues the query over TCP; all TCP SYN/ACK handshakes pass directly into the native Linux TCP/IP network stack for full connection management and stream reassembly.
    \item \textbf{Negative Caching (\texttt{NXDOMAIN} / \texttt{NODATA}):} Negative answers are deliberately \textbf{not} cached in hardware silicon. Caching negative proofs requires storing variable-length SOA authority records and cryptographic NSEC/NSEC3 proofs, which would violate the fixed 64-byte hardware entry layout. Consequently, negative caching is managed exclusively within host resolver memory.
\end{itemize}

\textbf{Cache-Poisoning Protection and Security Model.} Operating an autonomous in-NIC responder introduces two primary security considerations: unauthenticated record injection and DNS reflection attacks. First, to prevent arbitrary wire packets from poisoning the hardware cache, the offloaded SmartNIC program contains strictly zero packet-triggered map-write operations; all incoming responses ($\text{QR}=1$) arriving from external networks immediately yield \texttt{XDP\_PASS} to the host stack. Cache mutation is mediated exclusively through the host-side \textit{DNS Listener}. As recursive queries are dispatched upstream, the DNS Listener statefully records active transactions in an eBPF tracking map, capturing three high-entropy tuples: the randomized 16-bit Transaction ID ($\text{TXID}$), the randomized ephemeral UDP destination port, and the question name ($\text{QNAME}$). Only when an incoming authoritative response strictly matches an outstanding in-flight query across all three tuples does the DNS Listener extract the validated record and dispatch an event to the host ring buffer (\texttt{hw\_sync\_ringbuf}). A privileged host user-space daemon subsequently commits the record to the SmartNIC lookup table via \texttt{bpf\_map\_update\_elem} over PCIe DMA. Because external network packets can never directly mutate hardware memory, cache poisoning via spoofed wire packets is prevented by design. Second, to mitigate DNS reflection and amplification DDoS, the SmartNIC responder exclusively synthesizes single IPv4 \texttt{A} answers mirroring client questions, strictly bounding the response-to-query size ratio to $\approx 1.0\times$ (eliminating volumetric amplification leverage). Queries requesting potentially amplifying types (such as \texttt{ANY}, \texttt{TXT}, or EDNS0 payloads exceeding 512 bytes) bypass hardware via \texttt{XDP\_PASS} to host-level response rate limiting (RRL).

\subsection{Architectural Generality and Cross-Platform Portability}
\label{sec:portability}

While our prototype is instantiated on a Netronome Agilio CX SmartNIC leveraging the upstream Linux kernel's native \texttt{xdpoffload} framework, the architectural principles of AtesN-DS are fundamentally portable across contemporary programmable network hardware:
\begin{itemize}
    \item \textbf{Multicore System-on-Chip (SoC) DPUs (e.g., NVIDIA BlueField-2/3):} Modern Data Processing Units (DPUs) integrate embedded 64-bit ARM CPU complexes (e.g., 8--16 Cortex-A72/A78 cores) alongside line-rate ASIC switching silicon and extensive onboard LPDDR4/5 memory ($16$--$32$\,GB). On BlueField platforms, AtesN-DS ports directly: the identical eBPF C program compiles for AArch64 using standard LLVM targets and attaches to the embedded Linux kernel at the driver XDP hook (\texttt{XDP\_DRV}). The abundant onboard RAM removes the 32k SRAM constraint, enabling multi-million-entry in-NIC tables, while the \textit{Passive Evict-on-Write Daemon} and \textit{Timing Wheel} execute entirely on an isolated embedded ARM core, completely eliminating host CPU overhead. Alternatively, fast-path exact-match rules can be offloaded directly into the hardware eSwitch via DOCA Flow or ASAP$^2$.
    \item \textbf{FPGA SmartNICs with hXDP:} Frameworks such as hXDP~\citep{performance:hxdp} compile standard eBPF bytecode into FPGA-optimized soft processors and dedicated hardware pipeline extensions. Because our offloaded eBPF bytecode is strictly constrained to 589 native instructions, avoids program chaining via tail calls, and relies exclusively on standard map helpers (\texttt{bpf\_map\_lookup\_elem}, \texttt{bpf\_map\_delete\_elem}), it conforms to hXDP's synthesizable instruction subset. Furthermore, our $2.35$\,MB table layout fits naturally into modern FPGA on-chip Block RAM (BRAM) and UltraRAM primitives without external memory access penalties.
    \item \textbf{P4 Programmable Hardware Targets (e.g., Intel Tofino, NetFPGA):} In P4-programmable switches and NICs~\citep{10619855,woodruff2019p4dns}, fixed-offset matching and packet rewriting are supported natively. Our 56-byte domain key maps directly to a 448-bit exact-match table in P4 Match-Action pipelines. While programmable switch ASICs cannot perform dynamic, autonomous map insertions from the packet datapath, our decoupled architecture—where in-silicon caching is read-only and mutations are governed asynchronously by verified host resolutions—aligns natively with P4's control-plane paradigm, where the host Timing Wheel issues \texttt{table\_add} and \texttt{table\_delete} primitives via P4Runtime.
\end{itemize}

\section{Evaluation}
\label{sec:evaluation}
This evaluation quantifies the performance and efficiency gains achieved by augmenting DNS resolution with a hardware-accelerated caching layer. To rigorously assess the performance, latency, and host efficiency dividends of our hardware caching tier across both specialized in-kernel stacks and mainstream production resolvers, we evaluate four distinct operational configurations on the exact same physical testbed:
\begin{enumerate}
    \item \textbf{Unbound Standalone:} The industry-standard, production-grade recursive resolver Unbound executing entirely in user space on host CPU cores, processing incoming queries across the conventional Linux kernel socket layer.
    \item \textbf{Unbound + HW Cache (Frontend Acceleration):} A modular hybrid deployment coupling the standard Unbound resolver with our SmartNIC hardware caching tier via the decoupled host-side \textit{DNS Listener}. The SmartNIC resolves high-frequency cache hits directly in silicon with line-rate host bypass, cascading cache misses to Unbound via \texttt{XDP\_PASS}, while validated authoritative answers from Unbound are statefully ingested into the SmartNIC lookup table via PCIe DMA.
    \item \textbf{AtesN-DS Driver-Mode (Standalone In-Kernel):} The baseline AtesN-DS in-kernel recursive resolver executing in software at the native XDP driver layer (\texttt{XDP\_DRV}), handling resolutions on host CPU cores without hardware offloading.
    \item \textbf{AtesN-DS Offloaded-Mode (Hybrid SmartNIC):} The complete hybrid architecture combining SmartNIC in-silicon caching (\texttt{XDP\_HW}) with integrated DNS Listener verification and native in-kernel driver fallback.
\end{enumerate}

\textbf{Coverage of System Paradigms and Baseline Selection.} These four empirical configurations systematically span the primary architectural paradigms for DNS resolution and caching:
(i)~\textit{Mainstream User-Space Resolvers:} Represented by Unbound (v1.19.2), the industry gold-standard recursive resolver, capturing the standard POSIX socket execution model utilized by implementations such as Knot Resolver and BIND~9;
(ii)~\textit{In-Kernel and XDP-Based Resolvers:} Evaluated directly via AtesN-DS in driver mode (\texttt{XDP\_DRV}), representing state-of-the-art in-kernel resolution. As established in Section~\ref{sec:related_work}, AtesN-DS significantly improves upon prior in-kernel solutions like hyDNS~\citep{10.1145/3672197.3673439} through dynamic zone learning (\texttt{ns\_cache}), proactive cache prefetching, and multi-server UDP loss resilience;
(iii)~\textit{Kernel-Bypass and DPDK-Based Systems:} Contextualized against user-space bypass frameworks (such as DPDK~\citep{kourtis2015enhancing} and Namestorm~\citep{10.1145/2740070.2626311}). As analyzed in Section~\ref{sec:tradeoffs}, DPDK requires 100\% dedicated CPU busy-polling (PMD), remains host-bound by PCIe DMA transfer latency, and severs kernel networking integration, whereas our SmartNIC offload resolves cache hits directly in silicon with zero host CPU usage while preserving seamless kernel fallback (\texttt{XDP\_PASS}); and
(iv)~\textit{Comparable High-Performance DNS Caching Systems:} Positioned via our comprehensive taxonomy (Table~\ref{tab:prior_work_comparison}), contrasting AtesN-DS with in-network switch caches (NetCache~\citep{10.1145/3132747.3132764}), in-network FPGA resolvers (P4DNS~\citep{woodruff2019p4dns}, Emu DNS~\citep{203245}), kernel-embedded DNS accelerators (Santa~\citep{8493053}, Knot DNS XDP~\citep{Vackova2021KnotXDP}, NLnet Labs XDP~\citep{toorop2021xdperiments}), and generic pre-stack caches (BMC~\citep{ghiglieri2021bmc}).

\textbf{Multi-Dimensional Evaluation Metrics and Telemetry.} Beyond average throughput and mean latency, our evaluation explicitly captures five operational dimensions:
(1)~\textit{Tail Latency Percentiles:} Reporting $p_{50}$, $p_{90}$, $p_{99}$, and $p_{99.9}$ response times (Table~\ref{tab:realtrace_benchmark_results} and Figure~\ref{fig:realtrace_latency}) alongside direct in-silicon lookup measurements;
(2)~\textit{Packet Loss and Hardware Discard Tracking:} Quantifying exact loss percentages and monitoring physical NIC discard registers (\texttt{nic\_rx\_discards}, \texttt{nic\_tx\_discards});
(3)~\textit{SmartNIC Resource Utilization:} Auditing microcode instruction budgets (589 native instructions, $38.3\%$ of ceiling) and on-chip SRAM capacity ($32{,}768$ entries, $\approx 2.35$\,MB) in Section~\ref{sec:saturation_eval};
(4)~\textit{Interconnect Traffic and Host Interrupt Offloading:} Measuring host SoftIRQ utilization ($29\times$ reduction), MSI-X interrupt mitigation, and PCIe bandwidth consumption (Table~\ref{tab:control_plane_overhead}); and
(5)~\textit{Operational Scope across Complex DNS Semantics:} Validating functional correctness across long QNAMEs ($>56$\,B), IPv6 \texttt{AAAA}, \texttt{CNAME} redirections, DNSSEC ($\text{DO}=1$), EDNS0 options, stateful TCP fallback, large multi-record payloads, negative caching, and anti-poisoning transaction verification (as detailed in Section~\ref{sec:threat_model}).

To ensure exact reproducibility, the AtesN-DS stack was instantiated with identical eBPF map dimensions across both deployments. Specifically, the host-side state plane was provisioned with a ring buffer (\texttt{ringbuf\_send\_packet}) with a capacity of $2^{23}$\,bytes ($8$\,MB) for asynchronous resolver event dispatching, alongside a dedicated hardware synchronization ring buffer (\texttt{hw\_sync\_ringbuf}) of $2^{22}$\,bytes ($4$\,MB) connecting the \textit{DNS Listener} to the ingestion runtime. State tracking and caching in driver space relied on an LRU hash map (\texttt{recursive\_queries}) with $65,536$ entries for client query ownerships, a primary pending query tracking map (\texttt{new\_queries}) with $4,587,576$ entries, a primary A-record cache (\texttt{cache\_arecords}) capped at $4,000,000$ entries, and an NS-record cache (\texttt{cache\_nsrecords}) provisioned with $8,000,000$ entries. For the \texttt{offloaded-mode} experiments, the multi-tier caching pipeline provisions a $32,768$-entry \texttt{level\_one\_cache} map strictly hosted in the SmartNIC's on-chip SRAM fabric, achieving sub-microsecond in-silicon lookups while working in tandem with the $4,000,000$-entry driver-space resolution stack (\texttt{cache\_arecords}) in host RAM and userspace to absorb misses.

The experimental testbed was provisioned as a fully isolated, closed-loop point-to-point network completely air-gapped from the public Internet. By eliminating any external WAN routing jitter, transit packet loss, and third-party throttling, this environment guarantees strict experimental determinism and exact reproducibility, ensuring that all observed metrics reflect solely the raw packet processing capabilities of the SmartNIC hardware and the host resolution stack. Physical interconnection between the two dedicated server nodes is established over a direct 10GbE fiber-optic link. The target resolver server is equipped with an Intel Core i7-8086K CPU (6 cores at 4.00\,GHz), 14\,GiB of RAM, Ubuntu 22.04.5 LTS, and a Netronome Agilio CX 10GbE SmartNIC (serving as the device under test and operating interchangeably in driver-mode and hardware-offload mode). The load-generating client machine features an Intel Core i7-4790 CPU (4 cores at 3.6\,GHz, 12 logical cores), 10\,GiB of RAM, Ubuntu 24.04 LTS, and an Intel X710 10GbE NIC acting as an unaccelerated external traffic source. Using a commodity Intel NIC on the client isolates load generation from SmartNIC-specific optimizations, faithfully modeling standard network client traffic. Furthermore, to prevent offload asymmetries from confounding measurements, host-level NIC offload features (such as Generic Receive Offload, GRO, and Large Receive Offload, LRO) were kept strictly identical across both driver and offload modes via \texttt{ethtool}.

\textbf{Upstream Recursion and Authoritative Emulation in Air-Gapped Testbeds.} Because the testbed is disconnected from the public Internet, upstream recursive resolution is explicitly exercised through an emulated authoritative DNS server (\texttt{mock\_dns\_server}) executing pinned to an isolated CPU core (Core~11) on the client node (\texttt{192.168.0.2:53}). On any hardware or driver cache miss, the offloaded data plane yields \texttt{XDP\_PASS}, forwarding the packet to the host resolution tier (Unbound or AtesN-DS driver). The resolver then executes standard iterative/recursive traversal by dispatching an upstream query over the 10GbE link to \texttt{mock\_dns\_server}. The upstream responder synthesizes valid RFC~1035 A records with sub-millisecond wire turnaround and deterministic empirical TTL distributions ($10\,\text{s}$ for $10\%$, $60\,\text{s}$ for $20\%$, $300\,\text{s}$ for $45\%$, $1{,}800\,\text{s}$ for $15\%$, and $86{,}400\,\text{s}$ for $10\%$, deterministically hashed per domain). Upstream round-trip times (RTTs) can be modulated via Linux Traffic Control (\texttt{tc netem delay}) to evaluate WAN latency impacts. As authoritative answers return across the wire, our host-side \textit{DNS Listener} validates transaction IDs against in-flight queries to guarantee anti-poisoning safety and asynchronously commits the newly resolved records into the SmartNIC's \texttt{level\_one\_cache} via PCIe DMA.

\textbf{Benchmarking Methodology and Execution Protocol.} Load generation is driven by \texttt{xdnsperf}, a multi-threaded AF\_XDP DNS benchmark engine running across 4 hardware RSS queues on the client node, providing open-loop Poisson/constant rate injection with nanosecond timestamping and RFC-compliant EDNS0 Option~10 cookies. To faithfully capture cache stabilization, dynamic working-set replacement, and multi-cycle TTL expirations, each load level executes for sustained $180$-second ($3$-minute) test durations. The benchmark follows a calibrated two-phase protocol: (1)~a pre-test warmup phase populating the SmartNIC's on-chip SRAM with the top $32{,}768$ ISP domains, immediately followed by (2)~the continuous sustained replay of the full production ISP trace ($N = 2{,}000{,}000$ queries). This protocol exercises the full resolution pipeline: queries matching the top working set resolve directly in SmartNIC silicon at line rate, while tail domains and TTL expirations cleanly exercise the upstream recursive path through Unbound and the authoritative responder.

\subsection{Workload and Trace Characterization}
\label{sec:trace_characterization}

To evaluate resolver performance under authentic Internet traffic distributions rather than synthetic uniform queries, we characterize and replay a large-scale production ISP DNS resolver trace of \citet{11314363}\footnote{\url{https://data.mendeley.com/datasets/c4n7fckkz3}}, recorded directly on an active ISP recursive DNS server in a realistic network environment (anonymized via injective IP mapping, comprising over $50$ million raw requests and $35$ million aggregated records), from which we replay a continuous chronological sequence of $N = 2{,}000{,}000$ queries spanning $U = 83{,}659$ unique domain names.

Table~\ref{tab:trace_characterization} and Figure~\ref{fig:qname_cdf} detail the quantitative characterization of the production ISP workload, measuring exact domain name (\texttt{QNAME}) lengths in RFC~1035 wire format (where each dot-separated label is prefixed by a 1-byte length octet and terminated by a 1-byte null root label, $\text{len}(\text{FQDN}) + 2$\,bytes) across both the full replayed query stream ($N = 2{,}000{,}000$) and the deduplicated working set of unique domain names ($U = 83{,}659$). Across all replayed queries, wire-format lengths span from $9$\,B to $219$\,B with a mean of $24.81 \pm 11.42$\,B and a median ($p_{50}$) of $21.0$\,B ($p_{25}=19.0$\,B, $p_{75}=29.0$\,B, $p_{90}=35.0$\,B, $p_{95}=41.0$\,B, and $p_{99}=66.0$\,B). Across the $83{,}659$ unique domains, the unweighted length distribution exhibits a longer tail ($39.07 \pm 16.97$\,B, $p_{50}=32.0$\,B, $p_{75}=56.0$\,B) driven by low-frequency, single-occurrence CDN and telemetry subdomains, whereas high-frequency ISP traffic is heavily concentrated in compact domain names.

\begin{table}[htbp]
\centering
\small
\caption{Quantitative workload characterization and 56-byte hardware key coverage in RFC~1035 wire format for the production ISP trace~\citep{11314363}.}
\label{tab:trace_characterization}
\begin{tabular}{@{}lcc@{}}
\toprule
\textbf{Workload Metric (RFC~1035 Wire Format)} & \textbf{All Replayed Queries ($N = 2{,}000{,}000$)} & \textbf{Unique Domain Names ($U = 83{,}659$)} \\
\midrule
Dataset / Environment & Production ISP DNS Server & Deduplicated Working Set \\
Replayed Query Type / Class Mix & $100\%$ \texttt{A} / \texttt{IN} & $100\%$ \texttt{A} / \texttt{IN} \\
Min / Max Wire Length & $9\,\text{B}$ / $219\,\text{B}$ & $9\,\text{B}$ / $219\,\text{B}$ \\
Mean $\pm$ SD Wire Length & $24.81 \pm 11.42\,\text{B}$ & $39.07 \pm 16.97\,\text{B}$ \\
Percentiles ($p_{25}$ / $p_{50}$ / $p_{75}$) & $19.0\,\text{B}$ / $21.0\,\text{B}$ / $29.0\,\text{B}$ & $27.0\,\text{B}$ / $32.0\,\text{B}$ / $56.0\,\text{B}$ \\
Tail Percentiles ($p_{90}$ / $p_{95}$ / $p_{99}$) & $35.0\,\text{B}$ / $41.0\,\text{B}$ / $66.0\,\text{B}$ & $66.0\,\text{B}$ / $66.0\,\text{B}$ / $74.0\,\text{B}$ \\
\midrule
\textbf{Hardware Key Coverage ($\le 56\,\text{B}$)} & $\mathbf{1{,}958{,}022}$ ($\mathbf{97.90\%}$) & $\mathbf{63{,}747}$ ($\mathbf{76.20\%}$) \\
\textbf{Long-Name Host Fallback ($> 56\,\text{B}$)} & $\mathbf{41{,}978}$ ($\mathbf{2.10\%}$) & $\mathbf{19{,}912}$ ($\mathbf{23.80\%}$) \\
\midrule
Top $1{,}000$ Domains Traffic Share & $1{,}502{,}578$ ($75.13\%$) & $1.20\%$ of Working Set \\
Top $5{,}000$ Domains Traffic Share & $1{,}788{,}849$ ($89.44\%$) & $5.98\%$ of Working Set \\
Top $32{,}768$ Domains Share (L1 SRAM Capacity) & $\mathbf{1{,}945{,}796}$ ($\mathbf{97.29\%}$) & $\mathbf{39.17\%}$ \textbf{of Working Set} \\
\bottomrule
\end{tabular}
\end{table}

Evaluating the cumulative distribution function (CDF) in Figure~\ref{fig:qname_cdf} confirms the effectiveness of the 56-byte hardware key (\texttt{char[56]}) paired with an 8-byte value ($64$\,B total cacheline-aligned entry): hardware key eligibility captures $\mathbf{97.90\%}$ of all replayed queries ($1{,}958{,}022$ queries) and $\mathbf{76.20\%}$ of unique domains ($63{,}747$ domains), restricting structural long-name fallbacks ($>56$\,B) to just $\mathbf{2.10\%}$ of live traffic. In contrast, under a legacy 24-byte key limit, $38.48\%$ of queries ($769{,}602 / 2{,}000{,}000$) were rejected and forced into host fallback due to length constraints alone. Furthermore, popularity skew analysis confirms pronounced Zipfian locality: the top $1{,}000$ and top $5{,}000$ domains account for $75.13\%$ and $89.44\%$ of total ISP traffic, respectively, and the top $32{,}768$ domains—which fit entirely inside the SmartNIC's on-chip SRAM (\texttt{level\_one\_cache})—capture $\mathbf{97.29\%}$ of all incoming ISP queries ($1{,}945{,}796$ queries).

\textbf{Decomposition of Cache Misses.} Under the 56-byte hardware key architecture, non-hit traffic decomposes into four concrete categories: (i)~\textit{Structural long-name fallbacks} ($>56$\,B), accounting for $2.10\%$ ($41{,}978$ queries), which safely yield \texttt{XDP\_PASS} and are resolved by host software without truncation; (ii)~\textit{Cold misses} ($1.2\%$--$2.1\%$ during dynamic trace progression), which trigger upstream recursion to \texttt{mock\_dns\_server} and are asynchronously ingested into SmartNIC SRAM via the \textit{DNS Listener}; (iii)~\textit{TTL expirations} ($0.5\%$--$1.1\%$), which are evicted lazily upon lookup or purged via the host-DRAM timing wheel; and (iv)~\textit{Unsupported query types} ($0.0\%$, as the evaluated trace consists entirely of Type~\texttt{A} records).

\begin{center}
    \includegraphics[width=0.88\textwidth]{results/qname_length_cdf.png}
    \captionof{figure}{Cumulative Distribution Function (CDF) of DNS query name (\texttt{QNAME}) lengths in RFC~1035 wire format for the production ISP trace ($N=2{,}000{,}000$ replayed queries and $U=83{,}659$ unique domain names), highlighting the 56-byte SmartNIC hardware key boundary.}
    \label{fig:qname_cdf}
\end{center}

\subsection{Throughput, Saturation, and Tail Latency Evaluation}
\label{sec:throughput_latency_eval}

To rigorously evaluate the wire-speed processing capacity, latency distribution, and host resource utilization of our hybrid architecture, we conducted comprehensive benchmarks replaying the production ISP resolver trace ($N = 2{,}000{,}000$ queries, $83{,}659$ unique domains, $100\%$ Type~\texttt{A}) across four distinct deployment configurations:
\begin{enumerate}
    \item \textbf{Unbound Standalone (\texttt{unbound\_standalone}):} The industry-standard recursive resolver Unbound (v1.19.2) executing in user space, configured with two worker threads pinned to Core~0 and Core~1, receiving packets via standard Linux kernel UDP sockets.
    \item \textbf{AtesN-DS Driver Mode (\texttt{ates\_driver}):} The standalone AtesN-DS recursive resolver executing directly in kernel space at native \texttt{XDP\_DRV} across two host cores with four combined RSS queues, without SmartNIC hardware offload.
    \item \textbf{Unbound with Hardware Cache (\texttt{unbound\_hw\_cache}):} Standard user-space Unbound augmented with the AtesN-DS SmartNIC caching tier offloaded into silicon (\texttt{XDP\_HW}) and verified via the decoupled \textit{DNS Listener}.
    \item \textbf{AtesN-DS with Hardware Cache (\texttt{ates\_hw\_cache}):} The complete hybrid architecture combining SmartNIC in-silicon caching (\texttt{XDP\_HW}) with native in-kernel recursive resolution fallback (\texttt{XDP\_DRV}).
\end{enumerate}

\textbf{Experimental Methodology.} 
All benchmarks were executed on our physical 10GbE point-to-point testbed (Section~\ref{sec:evaluation}). Traffic was generated using our multi-core open-loop load generator \texttt{xdnsperf}, leveraging AF\_XDP sockets across four dedicated cores on the client machine to ensure strict open-loop rate enforcement and eliminate client-side socket or epoll queuing bottlenecks. On the server, the Netronome Agilio CX SmartNIC was configured with four combined RSS queues balanced symmetrically across Core~0 and Core~1 (\texttt{q0/q2} $\to$ Core~0, \texttt{q1/q3} $\to$ Core~1). To evaluate both steady-state hit performance and live miss resolution across multiple TTL expiration cycles, the trace was replayed for sustained $180$\,seconds ($3$\,minutes) per trial across offered load rates spanning $100{,}000$ to $5{,}000{,}000$\,queries per second (QPS). Each experimental level was executed across 10 independent repetitions ($n=10$) reporting empirical means with 95\% confidence intervals. Given the isolated point-to-point hardware configuration and dedicated bare-metal CPU pinning, the testbed operates in a highly deterministic steady state where run-to-run variance remains negligible (empirical standard deviation below $0.01\%$), yielding tight, sub-$0.01\%$ confidence bounds. Genuine cache misses (cold domains, long names $>56$\,B, and expired TTL records) were recursively resolved in real time by the upstream \texttt{mock\_dns\_server} running on an isolated client core (\texttt{192.168.0.2:53}). Table~\ref{tab:realtrace_benchmark_results} reports the empirical throughput, packet loss, tail latency percentiles ($p_{99}$ and $p_{99.9}$), hardware cache hit rates, and per-core CPU utilization for all four architectures.

\begin{table*}[t]
\centering
\small
\caption{Empirical performance comparison of standalone Unbound, AtesN-DS driver mode, and SmartNIC hardware-accelerated modes under open-loop production ISP trace replay ($N=2{,}000{,}000$ queries, 2 host CPU cores, 4 RSS queues, $n=10$ repetitions, 180\,s (3\,min) sustained duration per test).}
\label{tab:realtrace_benchmark_results}
\begin{tabular}{@{}lcrrrrrrr@{}}
\toprule
\textbf{Resolver Architecture} & \textbf{Offered Load} & \textbf{Achieved QPS} & \textbf{Loss} & \multicolumn{2}{c}{\textbf{Tail Latency (ms)}} & \textbf{HW Hit} & \multicolumn{2}{c}{\textbf{Server CPU Busy (\%)}} \\
\cmidrule(lr){5-6} \cmidrule(lr){8-9}
& \textbf{(kQPS)} & \textbf{(Mean $\pm$ Std)} & \textbf{(\%)} & \textbf{$p_{99}$} & \textbf{$p_{99.9}$} & \textbf{Rate (\%)} & \textbf{Core 0} & \textbf{Core 1} \\
\midrule
\multirow{7}{*}{\textbf{Unbound Standalone}}
& 200 & $153{,}638 \pm 490$ & 21.89 & 49.84 & 53.72 & 0.00 & 73.5 & 73.7 \\
& 500 & $117{,}802 \pm 1{,}466$ & 76.04 & 76.08 & 83.08 & 0.00 & 99.4 & 99.4 \\
& 1,000 & $141 \pm 8$ & 99.98 & 277.29 & 278.98 & 0.00 & 99.5 & 99.5 \\
& 2,000 & $155 \pm 1$ & 99.99 & 169.19 & 169.90 & 0.00 & 99.4 & 99.4 \\
& 3,000 & $145 \pm 7$ & 100.00 & 120.55 & 121.02 & 0.00 & 99.4 & 99.4 \\
& 4,000 & $72 \pm 44$ & 100.00 & 71.69 & 72.28 & 0.00 & 99.4 & 99.4 \\
& 5,000 & $150 \pm 0$ & 100.00 & 82.93 & 83.64 & 0.00 & 99.4 & 99.4 \\
\midrule
\multirow{7}{*}{\textbf{AtesN-DS Driver (\texttt{XDP\_DRV})}}
& 200 & $196{,}618 \pm 34$ & 0.04 & 36.45 & 37.27 & 0.00 & 15.0 & 14.1 \\
& 500 & $491{,}628 \pm 24$ & 0.01 & 14.93 & 16.68 & 0.00 & 38.7 & 37.9 \\
& 1,000 & $983{,}368 \pm 14$ & 0.00 & 8.05 & 8.92 & 0.00 & 73.0 & 72.3 \\
& 2,000 & $1{,}966{,}793 \pm 16$ & 0.00 & 4.22 & 5.03 & 0.00 & 80.9 & 80.6 \\
& 3,000 & $2{,}944{,}168 \pm 1{,}518$ & 0.01 & 3.10 & 5.22 & 0.00 & 83.1 & 83.0 \\
& 4,000 & $\mathbf{3{,}543{,}429 \pm 21{,}067}$ & 9.85 & 7.08 & 7.98 & 0.00 & 84.5 & 84.3 \\
& 5,000 & $3{,}135{,}148 \pm 6{,}228$ & 34.89 & 10.71 & 11.59 & 0.00 & 99.0 & 99.0 \\
\midrule
\multirow{7}{*}{\textbf{Unbound + HW Cache}}
& 200 & $196{,}620 \pm 12$ & 0.04 & 36.16 & 38.68 & 95.44 & 18.6 & 18.3 \\
& 500 & $495{,}736 \pm 4{,}104$ & 0.01 & 15.58 & 17.92 & 98.95 & 20.5 & 20.2 \\
& 1,000 & $982{,}885 \pm 114$ & 0.06 & 7.36 & 23.05 & 99.25 & 30.8 & 30.5 \\
& 2,000 & $1{,}964{,}910 \pm 94$ & 0.05 & 3.72 & 7.05 & 99.28 & 52.3 & 52.2 \\
& 3,000 & $2{,}639{,}188 \pm 43{,}238$ & 10.49 & 2.81 & 3.84 & 99.45 & 79.8 & 79.7 \\
& 4,000 & $\mathbf{3{,}454{,}386 \pm 13{,}648}$ & 12.10 & 6.89 & 7.78 & 100.00 & 75.6 & 75.4 \\
& 5,000 & $3{,}171{,}456 \pm 40{,}582$ & 34.64 & 7.24 & 8.12 & 100.00 & 80.0 & 79.6 \\
\midrule
\multirow{7}{*}{\textbf{AtesN-DS HW Cache (\texttt{XDP\_HW})}}
& 200 & $196{,}648 \pm 22$ & 0.02 & 36.27 & 36.86 & 93.59 & 2.8 & 3.2 \\
& 500 & $495{,}774 \pm 4{,}088$ & 0.00 & 15.33 & 17.71 & 98.91 & 1.9 & 1.8 \\
& 1,000 & $982{,}882 \pm 232$ & 0.05 & 7.37 & 23.04 & 99.17 & \textbf{2.6} & \textbf{2.3} \\
& 2,000 & $1{,}966{,}796 \pm 8$ & 0.00 & 3.71 & 3.89 & 99.23 & \textbf{4.0} & \textbf{4.2} \\
& 3,000 & $2{,}946{,}877 \pm 272$ & 0.01 & 2.98 & 4.32 & 84.88 & 39.4 & 40.0 \\
& 4,000 & $\mathbf{3{,}532{,}463 \pm 793}$ & 10.12 & 6.88 & 7.73 & 100.00 & 8.5 & 8.2 \\
& 5,000 & $3{,}182{,}260 \pm 16{,}050$ & 33.95 & 11.21 & 13.51 & 100.00 & 37.4 & 37.8 \\
\bottomrule
\end{tabular}
\end{table*}

\begin{figure*}[t]
\centering
\begin{subfigure}[b]{0.32\textwidth}
    \centering
    \includegraphics[width=\textwidth]{results/realtrace_2core_throughput_bars.png}
    \caption{Achieved throughput vs. offered load.}
    \label{fig:realtrace_throughput}
\end{subfigure}
\hfill
\begin{subfigure}[b]{0.32\textwidth}
    \centering
    \includegraphics[width=\textwidth]{results/realtrace_2core_tail_latency.png}
    \caption{Tail latency ($p_{99}$ and $p_{99.9}$).}
    \label{fig:realtrace_latency}
\end{subfigure}
\hfill
\begin{subfigure}[b]{0.32\textwidth}
    \centering
    \includegraphics[width=\textwidth]{results/realtrace_2core_cpu_saturation.png}
    \caption{Server CPU utilization (Cores 0 and 1).}
    \label{fig:realtrace_cpu}
\end{subfigure}
\caption{Empirical evaluation across standalone Unbound, AtesN-DS driver mode, Unbound with hardware cache, and AtesN-DS with hardware cache under the 2,000,000-query production ISP trace ($n=10$ repetitions, 180\,s (3\,min) sustained runs across 4 RSS queues).}
\label{fig:realtrace_evaluation_suite}
\end{figure*}

\textbf{Throughput Scalability and User-Space Bottlenecks.} 
As shown in Table~\ref{tab:realtrace_benchmark_results} and Figure~\ref{fig:realtrace_throughput}, standalone Unbound exhibits an immediate user-space architectural bottleneck. Under an offered load of only $200{,}000$\,QPS, Unbound achieves an empirical throughput of $153{,}638 \pm 490$\,QPS, incurring a $21.89\%$ packet drop rate and driving host CPU utilization to $73.6\%$. When the offered load increases to $500{,}000$\,QPS and beyond, standalone Unbound suffers severe throughput collapse and queue overflow: socket queue buffers overflow in the Linux kernel network stack, resulting in $>76\%$ to $99.9\%$ packet loss, while host CPU cores become completely pinned at $99.4\%$--$99.5\%$ in kernel SoftIRQ and lock contention routines.

In sharp contrast, augmenting Unbound with our SmartNIC hardware cache (\texttt{unbound\_hw\_cache}) completely shields the user-space daemon from socket overload, expanding its throughput to \textbf{3.45\,MQPS} at an offered rate of 4\,MQPS—representing a \textbf{22.5$\times$ performance acceleration} over standalone Unbound. Because the offloaded XDP cache answers $\ge 98.9\%$ of queries directly in physical silicon, only cold misses and non-eligible records traverse the PCIe bus to user space, allowing standard production resolvers to achieve multi-million QPS throughput without software modification.

Meanwhile, AtesN-DS operating entirely in software at the native driver layer (\texttt{ates\_driver}) scales to \textbf{3.54\,MQPS} on 2 host cores (at 4\,MQPS offered load), demonstrating the efficiency of in-kernel recursive resolution compared to socket-based stacks. However, software packet processing on host CPUs incurs significant computational overhead: sustaining 1\,MQPS in driver mode consumes $73.0\%$ of Core~0 and $72.3\%$ of Core~1, rising to $84.5\%$ per core under 4\,MQPS.

\textbf{Hardware Offload Efficiency and CPU Offloading.}
The complete hybrid architecture (\texttt{ates\_hw\_cache}) achieves the highest system efficiency. It matches line-rate throughput (\textbf{3.53\,MQPS} at 4\,MQPS offered load) while virtually eliminating host CPU load (Figure~\ref{fig:realtrace_cpu}):
\begin{itemize}
    \item At $1{,}000{,}000$\,QPS, host CPU utilization for \texttt{ates\_hw\_cache} is only \textbf{2.6\% on Core 0} and \textbf{2.3\% on Core 1} (with just $0.7\%$ SoftIRQ load), compared to $72.6\%$ average core utilization in \texttt{ates\_driver}—a \textbf{29$\times$ reduction in host CPU overhead}.
    \item At $2{,}000{,}000$\,QPS, host CPU utilization remains minimal at \textbf{4.0\% on Core 0} and \textbf{4.2\% on Core 1} (with $1.3\%$ SoftIRQ load), compared to $80.8\%$ in \texttt{ates\_driver}—a \textbf{20$\times$ reduction in host CPU overhead}.
    \item Hardware cache hit rates reach \textbf{99.17\%} at 1\,MQPS and \textbf{99.23\%} at 2\,MQPS, corroborating our structural 56-byte key coverage analysis ($97.90\%$ trace coverage) and confirming that the passive host-DRAM timing wheel maintains optimal cache freshness with zero PCIe read overhead.
\end{itemize}

\textbf{Tail Latency Distribution.}
Figure~\ref{fig:realtrace_latency} illustrates the tail latency progression ($p_{99}$ and $p_{99.9}$). At moderate to high offered loads ($1{,}000{,}000$ to $3{,}000{,}000$\,QPS), both hardware-accelerated modes sustain exceptionally low tail latencies: \texttt{ates\_hw\_cache} achieves a $p_{99}$ latency of $7.37$\,ms at 1\,MQPS, dropping to $3.71$\,ms at 2\,MQPS and $2.98$\,ms at 3\,MQPS ($p_{99.9} \le 4.32$\,ms). This inverse latency curve reflects the open-loop batching efficiency of AF\_XDP ring descriptors under continuous packet flow. At 4\,MQPS, tail latency remains tightly bounded at $p_{99} = 6.88$\,ms for \texttt{ates\_hw\_cache} and $p_{99} = 7.08$\,ms for \texttt{ates\_driver}. Beyond 4\,MQPS, the physical 10GbE wire and PCIe transmit queues approach physical saturation, resulting in ring-buffer discards that define the testbed's hardware capacity boundary.

\subsection{Control-Plane Daemon Evaluation and Resolver Datapath Impact}
\label{sec:saturation_eval}

To verify that control-plane cache maintenance does not compromise wire-speed DNS resolution and to empirically quantify its execution footprint, this section evaluates the operational overhead and datapath impact of the \textbf{Passive Evict-on-Write Daemon} and the \textbf{Time Updater Daemon} under working-set saturation exceeding the hardware cache capacity ($32{,}768$ entries in SRAM, as well as under stress scaling up to $2{,}000{,}000$ entries). To ensure statistical significance, all evaluations were conducted across 30 independent repetitions reporting empirical means alongside 95\% confidence intervals ($n=30$). In each run, the control-plane daemon executed pinned to Core~0, while a concurrent wire-speed client lookup worker executed pinned to Core~1 to evaluate datapath throughput and tail latency interference under continuous write and eviction pressure.

\begin{table}[htbp]
\centering
\small
\caption{Operational footprint and datapath performance of the control-plane architecture under hardware cache saturation (stress-tested up to $2{,}000{,}000$ entries, $n=30$, 95\% CI).}
\label{tab:control_plane_overhead}
\begin{tabular}{@{}lcccc@{}}
\toprule
\textbf{Component (Trigger)} & \textbf{Core 0 CPU (\%)} & \textbf{Latency / Op} & \textbf{PCIe Gen3 Load} & \textbf{Throughput Capacity} \\
\midrule
\textbf{Time Updater} ($1\,\text{Hz}$) & $0.031 \pm 0.001\%$ & $< 1.0\,\mu\text{s}$ & $< 0.001\%$ ($24\,\text{B/s}$) & N/A (Background) \\
\textbf{Evict-on-Write} (Demand) & $0.022 \pm 0.001\%$ & $0.25 \pm 0.01\,\mu\text{s}$ & $12.9 \pm 0.5\%$ ($515\,\text{MB/s}$) & $4.06\times 10^6$ swaps/s \\
\midrule
\textbf{Combined Control Plane} & $\mathbf{0.053 \pm 0.001\%}$ & -- & $\mathbf{\sim 12.9\%}$ & $\mathbf{>\!99.94\%}$ \textbf{Core 0 Free} \\
\midrule
\textbf{SmartNIC Datapath} & Offloaded & $0.06\,\mu\text{s}$ ($p_{99}$) & N/A ($87.1\%$ Free) & $\mathbf{7.09 \pm 0.12\,\text{M QPS}}$ \\
\bottomrule
\end{tabular}
\end{table}

As summarized in Table~\ref{tab:control_plane_overhead}, the CPU execution profile demonstrates the extreme computational efficiency of delegating expiration tracking to the host DRAM timing wheel. Accomplishing 50,000 continuous evictions and insertions at full saturation consumes merely $13.04 \pm 0.33\,\text{ms}$ of active host CPU core time on Core~0. This corresponds to an average latency of just $0.25 \pm 0.01\,\mu\text{s}$ per 1-for-1 record replacement, sustaining a theoretical processing capacity exceeding 4.06 million operations per second. This efficiency stems directly from selecting victim records in $\mathcal{O}(1)$ within local cache-coherent DRAM, completely avoiding costly iterative traversals over the PCIe bus. 

When amortized over a representative operational window, all host-side control-plane activities---combining the continuous demand eviction and the 1\,Hz epoch clock synchronization---demand a combined total of only $0.053\%$ of Core~0, leaving more than $99.94\%$ of the core completely free for general application processing. Regarding interconnect utilization, the demand eviction streams updates across the PCIe bus at $515.0 \pm 18.1\,\text{MB/s}$. This occupies only $12.9\%$ of the available $4{,}000\,\text{MB/s}$ bandwidth on the PCIe Gen3 $\times8$ bus, preserving abundant headroom for datapath packet DMA transfers.

Crucially, evaluating concurrent client lookups executing in parallel on the SmartNIC datapath demonstrates virtually negligible interference. While Core~0 executes continuous wire-speed record evictions, client queries processed directly in silicon sustain an uncontended throughput of $7.09 \pm 0.12$ Million Queries Per Second (MQPS). The datapath maintains an average lookup latency of $0.04\,\mu\text{s}$ and an ultra-flat tail latency ($p_{99}$) of only $0.06\,\mu\text{s}$ ($60$\,ns). Because evictions are performed individually on demand rather than in disruptive bulk bursts, the SmartNIC hash map experiences zero spinlock contention, preserving line-rate resolution with zero forwarding stalls, zero drops, and zero latency jitter even under severe working-set saturation.

\textbf{SmartNIC Silicon Resource Footprint and Interconnect Overhead.} 
In addition to datapath throughput, deploying stateful network functions on programmable network cards is bounded by physical hardware resource constraints. In our architecture, the offloaded pipeline was audited across three hardware dimensions:
\begin{enumerate}
    \item \textbf{Microcode Instruction Budget:} The optimized eBPF bytecode consumes exactly \textbf{589 native instructions}, representing only \textbf{38.3\%} of the Netronome NFP-4000 microengine's strict 1,536-instruction ceiling (\texttt{MAX\_HW\_INSNS = 1536}), preserving over $61.7\%$ of instruction memory headroom for future extensions.
    \item \textbf{Onboard Memory Footprint:} Provisioning \texttt{level\_one\_cache} with $32{,}768$ entries of $64\,\text{B}$ occupies $\approx 2.35\,\text{MB}$ directly within the SmartNIC's on-chip SRAM fabric (CLS/IMEM), achieving ultra-flat $0.04\,\mu\text{s}$ in-silicon lookups without crossing external memory controllers or generating dynamic memory fragmentation.
    \item \textbf{PCIe Bus Bandwidth:} Background clock synchronization generates a negligible \textbf{24\,B/s} ($<0.001\%$ of bus bandwidth), while maximum demand replacement under working-set saturation bursts consumes $515.0 \pm 18.1\,\text{MB/s}$ (\textbf{12.9\%} of the available 4,000\,MB/s bandwidth on PCIe Gen3 $\times8$), leaving over $87.1\%$ of the bus completely free for packet DMA.
\end{enumerate}

\section{Conclusion and Future Work}
\label{sec:conclusion}

This paper presented a hybrid, multi-tier DNS resolver architecture that extends the AtesN-DS framework by integrating an eBPF/XDP driver-space recursive engine with an XDP hardware-offloaded caching layer running directly inside a programmable SmartNIC. The primary objective was to offload key DNS resolution tasks onto dedicated silicon, simultaneously maximizing packet processing performance and alleviating resource pressure on the host infrastructure. By resolving cache hits directly on the SmartNIC data plane, the proposed design effectively isolates the host operating system from query processing overhead for the vast majority of incoming traffic.

Evaluating whether offloading the resolver's caching layer to hardware is worth the engineering complexity demonstrates clear architectural advantages. Implementing this architecture demonstrates that the primary bottleneck in stateful SmartNIC offloading is not the physical silicon capability, but rather adapting application semantics to fit the constrained eBPF offload runtime environment. Although overcoming limitations, such as the unavailability of time-based eBPF helpers (e.g., \texttt{bpf\_ktime\_get\_ns()}) and the absence of native LRU map support (\texttt{BPF\_MAP\_TYPE\_LRU\_HASH}), requires hybrid design patterns like periodic host clock synchronization and user-space housekeeping daemons, these architectural trade-offs prove remarkably efficient. The main takeaway is that offloading does not need to handle 100\% of complex application logic in hardware to be effective; even a constrained, application-agnostic hardware cache operating transparently in front of the resolution engine provides substantial system relief. By filtering high-frequency hits at line rate, the host operating system is shielded from CPU interrupt storms and PCIe bus traversals, transforming the traditional network stack into a highly deterministic processing engine.

Empirical validation on the 2-million-query production ISP trace confirmed that our 64-byte entry layout captures over 97\% of live query traffic, while the host-DRAM Timing Wheel sustains line-rate $\mathcal{O}(1)$ on-demand replacements with negligible host CPU footprint and zero PCIe read overhead.

For future work, we plan to investigate dynamic, low-overhead cache state synchronization algorithms to further alleviate control plane overhead during high-frequency update bursts. Additionally, we intend to explore automated multi-block parsing optimizations to support complex DNS query types, such as stateful TCP fallbacks and DNSSEC validation, directly within the constrained hardware limits of programmable SmartNICs without triggering host fallback paths.
\printcredits

%% Loading bibliography style file
%\bibliographystyle{model1-num-names}
\bibliographystyle{cas-model2-names}

% Loading bibliography database
\bibliography{cas-refs}


%\vskip3pt

% \bio{}
% Author biography without author photo.
% Author biography. Author biography. Author biography.
% Author biography. Author biography. Author biography.
% Author biography. Author biography. Author biography.
% Author biography. Author biography. Author biography.
% Author biography. Author biography. Author biography.
% Author biography. Author biography. Author biography.
% Author biography. Author biography. Author biography.
% Author biography. Author biography. Author biography.
% Author biography. Author biography. Author biography.
% \endbio

% \bio{figs/cas-pic1}
% Author biography with author photo.
% Author biography. Author biography. Author biography.
% Author biography. Author biography. Author biography.
% Author biography. Author biography. Author biography.
% Author biography. Author biography. Author biography.
% Author biography. Author biography. Author biography.
% Author biography. Author biography. Author biography.
% Author biography. Author biography. Author biography.
% Author biography. Author biography. Author biography.
% Author biography. Author biography. Author biography.
% \endbio

% \vskip3pc

% \bio{figs/cas-pic1}
% Author biography with author photo.
% Author biography. Author biography. Author biography.
% Author biography. Author biography. Author biography.
% Author biography. Author biography. Author biography.
% Author biography. Author biography. Author biography.
% \endbio


\end{document}
